\chapter{Teoremas y Demostraciones Auxiliares}
\label{ch:appendix}

\providecommand{\Q}{\mathbb{Q}}
\providecommand{\GL}{\operatorname{GL}}
\providecommand{\SO}{\operatorname{SO}}
\providecommand{\Orth}{\operatorname{O}}
\providecommand{\SU}{\operatorname{SU}}
\providecommand{\Sym}{\operatorname{Sym}}
\providecommand{\rank}{\operatorname{rank}}
\providecommand{\tr}{\operatorname{tr}}
\providecommand{\Img}{\operatorname{Im}}
\providecommand{\Ric}{\operatorname{Ric}}
\providecommand{\Vol}{\operatorname{Vol}}
\providecommand{\diag}{\operatorname{diag}}
\providecommand{\II}{\mathrm{II}}
\providecommand{\dvol}{\,dA}

Este ap\'endice re\'une teoremas cl\'asicos de geometr\'ia diferencial,
\'algebra lineal, teor\'ia de grupos, teor\'ia de la medida y an\'alisis
funcional, junto con sus demostraciones. Aunque estos resultados son
fundamentales para el rigor matem\'atico del texto principal, no son centrales
para la teor\'ia del aprendizaje profundo geom\'etrico desarrollada en los
cap\'itulos principales, y agruparlos aqu\'i mantiene la exposici\'on
enfocada. Cada secci\'on presenta un teorema con los preliminares que requiere,
un enunciado formal y una demostraci\'on completa o un esquema de demostraci\'on
con referencias a las fuentes originales. La maquinaria de teor\'ia de la medida
y an\'alisis funcional de las dos \'ultimas secciones sustenta el teorema de
aproximaci\'on universal, cuyo enunciado permanece en el \Cref{ch:foundations} y
que remite a este ap\'endice para sus prerrequisitos m\'as pesados.

%=======================================================================
\section{El Teorema de Gauss--Bonnet}
\label{sec:app:gauss-bonnet}
%=======================================================================

El teorema de Gauss--Bonnet es uno de los resultados m\'as profundos de la
geometr\'ia diferencial, pues establece una conexi\'on fundamental entre la
geometr\'ia local (la curvatura) y la topolog\'ia global (la caracter\'istica de
Euler) de una superficie. Demuestra que ciertas cantidades geom\'etricas est\'an
determinadas \'unicamente por la topolog\'ia de la superficie.

\subsection{Preliminares}
\label{subsec:app:gauss-bonnet-prelim}

Antes de enunciar el teorema, introducimos las nociones necesarias.

\begin{definition}[Curvatura gaussiana]
\label{def:app:gaussian-curvature}
Sea $\mathcal{M}$ una superficie regular inmersa en $\R^3$. En cada punto
$p \in \mathcal{M}$ consideramos las \emph{curvaturas principales} $\kappa_1$ y
$\kappa_2$, los valores propios de la segunda forma fundamental. La
\emph{curvatura gaussiana} $K \colon \mathcal{M} \to \R$ es el producto de las
curvaturas principales:
\begin{equation}
\label{eq:app:gaussian-curvature}
K = \kappa_1 \kappa_2 .
\end{equation}
\end{definition}

\begin{remark}
El \emph{Theorema Egregium} de Gauss establece que la curvatura gaussiana $K$ es
una \emph{invariante isom\'etrica}: depende solo de la primera forma fundamental
(la m\'etrica) y no de c\'omo la superficie est\'a inmersa en $\R^3$. As\'i, $K$
es una propiedad intr\'inseca de la superficie.
\end{remark}

\begin{definition}[Curvatura geod\'esica]
\label{def:app:geodesic-curvature}
Sea $\gamma \colon [a,b] \to \mathcal{M}$ una curva regular parametrizada por
longitud de arco sobre una superficie orientada $\mathcal{M}$. Sea $\mathbf{n}$
la normal unitaria a la superficie y $\mathbf{T} = \gamma'$ la tangente unitaria
a la curva. La \emph{curvatura geod\'esica} $\kappa_g$ de $\gamma$ es
\begin{equation}
\label{eq:app:geodesic-curvature}
\kappa_g = \inner{\gamma''}{\mathbf{n} \times \mathbf{T}} ,
\end{equation}
donde $\mathbf{n} \times \mathbf{T}$ es la normal a la curva que yace en el plano
tangente a la superficie.
\end{definition}

\begin{remark}
La curvatura geod\'esica $\kappa_g$ mide cu\'anto se desv\'ia una curva de ser una
geod\'esica dentro de la superficie. Una curva es una geod\'esica si y solo si
$\kappa_g = 0$ en todos sus puntos. Intuitivamente, $\kappa_g$ captura la
componente de la curvatura que es ``visible'' desde dentro de la superficie.
\end{remark}

\begin{definition}[Caracter\'istica de Euler]
\label{def:app:euler-characteristic}
Sea $\mathcal{M}$ una superficie compacta. La \emph{caracter\'istica de Euler}
$\chi(\mathcal{M})$ es un invariante topol\'ogico definido de las siguientes
maneras equivalentes:
\begin{enumerate}
    \item \textbf{Triangulaci\'on.} Si $\mathcal{M}$ admite una triangulaci\'on
          con $V$ v\'ertices, $E$ aristas y $F$ caras, entonces
          \begin{equation}
          \label{eq:app:euler-triangulation}
          \chi(\mathcal{M}) = V - E + F .
          \end{equation}
    \item \textbf{G\'enero.} Para una superficie orientable de g\'enero $g$
          (n\'umero de ``asas''),
          \begin{equation}
          \label{eq:app:euler-genus}
          \chi(\mathcal{M}) = 2 - 2g .
          \end{equation}
\end{enumerate}
\end{definition}

\begin{example}
\label{ex:app:euler-surfaces}
Algunos valores de la caracter\'istica de Euler para superficies conocidas:
\begin{itemize}
    \item Esfera $S^2$: $\chi(S^2) = 2$ (g\'enero $g = 0$).
    \item Toro $T^2$: $\chi(T^2) = 0$ (g\'enero $g = 1$).
    \item Doble toro: $\chi = -2$ (g\'enero $g = 2$).
    \item Disco: $\chi = 1$ (una superficie con borde).
\end{itemize}
\end{example}

\subsection{Enunciado del teorema}
\label{subsec:app:gauss-bonnet-statement}

Presentamos primero la versi\'on general para superficies con borde y luego
derivamos casos particulares importantes.

\begin{theorem}[Gauss--Bonnet para superficies con borde]
\label{thm:app:gauss-bonnet-general}
Sea $\mathcal{M}$ una superficie riemanniana compacta orientada con borde
$\partial \mathcal{M}$ suave a trozos. Supongamos que $\partial \mathcal{M}$
consta de curvas regulares unidas en v\'ertices donde forman \'angulos exteriores
$\alpha_1, \ldots, \alpha_n$. Entonces
\begin{equation}
\label{eq:app:gauss-bonnet-general}
\int_{\mathcal{M}} K \, dA + \int_{\partial \mathcal{M}} \kappa_g \, ds
+ \sum_{i=1}^n \alpha_i = 2\pi \chi(\mathcal{M}) ,
\end{equation}
donde $K$ es la curvatura gaussiana, $dA$ el elemento de \'area inducido por la
m\'etrica riemanniana, $\kappa_g$ la curvatura geod\'esica del borde, $ds$ el
elemento de longitud de arco y $\alpha_i$ el \'angulo exterior en el $i$-\'esimo
v\'ertice del borde.
\end{theorem}

\begin{corollary}[Gauss--Bonnet para superficies cerradas]
\label{cor:app:gauss-bonnet-closed}
Si $\mathcal{M}$ es una superficie compacta sin borde (cerrada), entonces
\begin{equation}
\label{eq:app:gauss-bonnet-closed}
\int_{\mathcal{M}} K \, dA = 2\pi \chi(\mathcal{M}) .
\end{equation}
\end{corollary}

\begin{proof}
Si $\partial \mathcal{M} = \emptyset$, la integral sobre el borde y la suma de
\'angulos exteriores se anulan, y el resultado se sigue directamente del
\Cref{thm:app:gauss-bonnet-general}.
\end{proof}

\begin{corollary}[Suma de \'angulos en tri\'angulos geod\'esicos]
\label{cor:app:gauss-bonnet-triangle}
Sea $T \subset \mathcal{M}$ un tri\'angulo geod\'esico (una regi\'on limitada por
tres geod\'esicas) en una superficie riemanniana, con \'angulos interiores
$\alpha$, $\beta$, $\gamma$. Entonces
\begin{equation}
\label{eq:app:gauss-bonnet-triangle}
\alpha + \beta + \gamma = \pi + \int_T K \, dA .
\end{equation}
\end{corollary}

\begin{proof}
Aplicamos el \Cref{thm:app:gauss-bonnet-general} al tri\'angulo $T$:
\begin{enumerate}
    \item Como $T$ es homeomorfo a un disco, $\chi(T) = 1$.
    \item Como los lados son geod\'esicas, $\kappa_g = 0$ en todo el borde, de
          modo que $\int_{\partial T} \kappa_g \, ds = 0$.
    \item Los \'angulos exteriores en los v\'ertices satisfacen
          $\alpha_i^{\text{ext}} = \pi - \alpha_i^{\text{int}}$, donde
          $\alpha_i^{\text{int}}$ es el \'angulo interior correspondiente.
\end{enumerate}
El teorema da, por tanto,
\[
\int_T K \, dA + 0 + \sum_{i=1}^3 (\pi - \theta_i) = 2\pi ,
\]
donde $\theta_1, \theta_2, \theta_3$ son los \'angulos interiores
$\alpha, \beta, \gamma$. Simplificando,
\[
\int_T K \, dA + 3\pi - (\alpha + \beta + \gamma) = 2\pi ,
\qquad\text{de donde}\qquad
\alpha + \beta + \gamma = \pi + \int_T K \, dA . \qedhere
\]
\end{proof}

\begin{remark}[Interpretaci\'on geom\'etrica]
El \Cref{cor:app:gauss-bonnet-triangle} tiene implicaciones profundas:
\begin{itemize}
    \item \textbf{Geometr\'ia euclidiana} ($K = 0$):
          $\alpha + \beta + \gamma = \pi$ (el resultado cl\'asico de Euclides).
    \item \textbf{Geometr\'ia esf\'erica} ($K > 0$):
          $\alpha + \beta + \gamma > \pi$ (exceso angular proporcional al
          \'area).
    \item \textbf{Geometr\'ia hiperb\'olica} ($K < 0$):
          $\alpha + \beta + \gamma < \pi$ (defecto angular proporcional al
          \'area).
\end{itemize}
Esto muestra que el quinto postulado de Euclides (sobre las paralelas) es
equivalente a la condici\'on $K = 0$.
\end{remark}

\subsection{Demostraci\'on del teorema}
\label{subsec:app:gauss-bonnet-proof}

Presentamos un esquema de demostraci\'on siguiendo el enfoque
cl\'asico~\citep{do1976differential,lee2018introduction}.

\begin{lemma}[F\'ormula para pol\'igonos geod\'esicos]
\label{lem:app:geodesic-polygon}
Sea $P$ un pol\'igono geod\'esico simple en una superficie $\mathcal{M}$ con $n$
v\'ertices y \'angulos interiores $\theta_1, \ldots, \theta_n$. Entonces
\begin{equation}
\label{eq:app:geodesic-polygon}
\sum_{i=1}^n \theta_i = (n-2)\pi + \int_P K \, dA .
\end{equation}
\end{lemma}

\begin{proof}
Triangulamos el pol\'igono $P$ en $n-2$ tri\'angulos geod\'esicos
$T_1, \ldots, T_{n-2}$ (siempre posible para pol\'igonos simples). Aplicando el
\Cref{cor:app:gauss-bonnet-triangle} a cada tri\'angulo y sumando,
\[
\sum_{j=1}^{n-2} (\alpha_j + \beta_j + \gamma_j)
= (n-2)\pi + \sum_{j=1}^{n-2} \int_{T_j} K \, dA .
\]
En el lado izquierdo, los \'angulos en v\'ertices interiores suman $2\pi$ cada uno
(\'angulo completo), mientras que los \'angulos en los v\'ertices del borde son
exactamente los $\theta_i$ del pol\'igono original. En el lado derecho, la suma de
integrales es la integral sobre todo $P$:
\[
\sum_{i=1}^n \theta_i + 2\pi \cdot (\text{v\'ertices interiores})
= (n-2)\pi + \int_P K \, dA .
\]
Un conteo cuidadoso de la triangulaci\'on (usando la f\'ormula de Euler para
grafos planares) muestra que la contribuci\'on de los v\'ertices interiores se
cancela, dando el resultado deseado.
\end{proof}

\begin{proof}[Esquema de demostraci\'on del \Cref{thm:app:gauss-bonnet-general}]
La demostraci\'on procede en varios pasos.

\textbf{Paso 1: Triangulaci\'on de la superficie.}
Por el \Cref{thm:app:rado}, toda superficie compacta admite una triangulaci\'on.
Refinando si es necesario, obtenemos una triangulaci\'on geod\'esica
$\mathcal{T} = \{T_1, \ldots, T_F\}$ en la que cada tri\'angulo tiene lados
geod\'esicos.

\textbf{Paso 2: Aplicaci\'on del lema a cada tri\'angulo.}
Aplicamos el \Cref{lem:app:geodesic-polygon} (caso $n = 3$) a cada tri\'angulo:
\[
\alpha_j + \beta_j + \gamma_j = \pi + \int_{T_j} K \, dA ,
\qquad j = 1, \ldots, F .
\]

\textbf{Paso 3: Suma sobre todos los tri\'angulos.}
Sumando sobre todos los tri\'angulos,
\[
\sum_{j=1}^F (\alpha_j + \beta_j + \gamma_j) = F\pi + \int_{\mathcal{M}} K \, dA .
\]

\textbf{Paso 4: An\'alisis de la suma de \'angulos.}
La suma de \'angulos del lado izquierdo se descompone seg\'un la ubicaci\'on de
los v\'ertices:
\begin{itemize}
    \item \textit{V\'ertices interiores}: cada uno contribuye $2\pi$ (\'angulo
          completo alrededor del punto). Si hay $V_{\text{int}}$ v\'ertices
          interiores, contribuyen $2\pi V_{\text{int}}$.
    \item \textit{V\'ertices del borde suaves}: cada uno en una parte suave del
          borde contribuye $\pi$ (medio \'angulo).
    \item \textit{V\'ertices del borde (esquinas)}: cada esquina con \'angulo
          interior $\theta_i$ contribuye exactamente $\theta_i = \pi - \alpha_i$.
\end{itemize}

\textbf{Paso 5: Aplicaci\'on de la f\'ormula de Euler.}
Sean $V$, $E$, $F$ el n\'umero total de v\'ertices, aristas y caras en la
triangulaci\'on. Por la f\'ormula de Euler, $V - E + F = \chi(\mathcal{M})$.
Adem\'as, contando aristas: cada tri\'angulo tiene tres aristas, cada arista
interior es compartida por dos tri\'angulos y las aristas del borde pertenecen a
un solo tri\'angulo, lo que da relaciones entre $V$, $E$, $F$ y las cantidades del
borde.

\textbf{Paso 6: Combinaci\'on de t\'erminos.}
Combinando las identidades de los pasos 3--5 con \'algebra cuidadosa, los
t\'erminos se reorganizan para dar
\[
\int_{\mathcal{M}} K \, dA + \int_{\partial \mathcal{M}} \kappa_g \, ds
+ \sum_{i=1}^n \alpha_i = 2\pi \chi(\mathcal{M}) .
\]
La integral de curvatura geod\'esica aparece naturalmente al pasar al l\'imite
cuando la triangulaci\'on se refina cerca del borde.
\end{proof}

\begin{remark}
Una demostraci\'on completa y rigurosa puede encontrarse en
\citet[Cap\'itulo~4]{do1976differential} o en
\citet[Cap\'itulo~9]{lee2018introduction}. La versi\'on para variedades
riemannianas abstractas (no necesariamente inmersas en $\R^3$) requiere conexiones
y transporte paralelo.
\end{remark}

\subsection{Generalizaciones}
\label{subsec:app:gauss-bonnet-generalizations}

El teorema de Gauss--Bonnet admite varias generalizaciones importantes:
\begin{enumerate}
    \item \textbf{Teorema de Gauss--Bonnet--Chern.} Para variedades riemannianas
          de dimensi\'on par $2n$, el teorema se generaliza usando la forma de
          Pfaffian de la curvatura~\citep{chern1944simple}.
    \item \textbf{Teorema del \'indice de Atiyah--Singer.} Una generalizaci\'on
          profunda que conecta invariantes anal\'iticos (el \'indice de operadores
          el\'ipticos) con invariantes topol\'ogicos.
    \item \textbf{Versi\'on discreta.} Para mallas poligonales existe una
          versi\'on discreta del teorema, \'util en geometr\'ia computacional y
          procesamiento de geometr\'ia.
\end{enumerate}

%=======================================================================
\section{El Teorema de Sard}
\label{sec:app:sard}
%=======================================================================

El teorema de Sard es una herramienta fundamental del an\'alisis en variedades
que permite argumentos de ``posici\'on general'': garantiza que, en un sentido
preciso, la mayor\'ia de los puntos del codominio de una aplicaci\'on suave son
valores regulares. Este resultado es esencial en la demostraci\'on del teorema de
inmersi\'on de Whitney.

\begin{definition}[Punto cr\'itico y valor cr\'itico]
\label{def:app:critical-point}
Sea $F \colon \mathcal{M} \to \mathcal{N}$ una aplicaci\'on suave entre variedades
diferenciables. Un punto $p \in \mathcal{M}$ es un \emph{punto cr\'itico} de $F$
si el diferencial $dF_p \colon T_p\mathcal{M} \to T_{F(p)}\mathcal{N}$ no es
sobreyectivo. Un punto $y \in \mathcal{N}$ es un \emph{valor cr\'itico} si es
imagen de alg\'un punto cr\'itico, es decir, si $y = F(p)$ para alg\'un punto
cr\'itico $p$.
\end{definition}

\begin{definition}[Valor regular]
\label{def:app:regular-value}
Un punto $y \in \mathcal{N}$ es un \emph{valor regular} de
$F \colon \mathcal{M} \to \mathcal{N}$ si no es valor cr\'itico, es decir, si para
todo $p \in F^{-1}(y)$ el diferencial $dF_p$ es sobreyectivo. En particular, si
$y \notin F(\mathcal{M})$, entonces $y$ es trivialmente un valor regular.
\end{definition}

\begin{theorem}[Teorema de Sard]
\label{thm:app:sard}
Sea $F \colon \mathcal{M} \to \mathcal{N}$ una aplicaci\'on suave entre variedades
diferenciables. Entonces el conjunto de valores cr\'iticos de $F$ tiene medida de
Lebesgue nula en $\mathcal{N}$. Equivalentemente, el conjunto de valores regulares
es denso en $\mathcal{N}$.
\end{theorem}

La demostraci\'on completa, que procede por inducci\'on sobre la dimensi\'on y usa
estimaciones de medida basadas en el desarrollo de Taylor de $F$, puede
encontrarse en \citet[Cap\'itulo~6, Teorema~6.10]{lee2012introduction} o en
\citet[Cap\'itulo~3]{milnor1965topology}.

\begin{remark}
La potencia del teorema de Sard reside en que permite concluir la existencia de
objetos con buenas propiedades (valores regulares) sin construirlos
expl\'icitamente: basta mostrar que los ``malos'' forman un conjunto de medida
nula. En la demostraci\'on del teorema de Whitney, este argumento garantiza la
existencia de direcciones de proyecci\'on que preservan el embebimiento.
\end{remark}

%=======================================================================
\section{El Teorema del Nivel Regular}
\label{sec:app:regular-level}
%=======================================================================

El teorema del nivel regular conecta la teor\'ia de aplicaciones suaves entre
variedades con la estructura de subvariedad de sus conjuntos de nivel. Es una
herramienta fundamental de la geometr\'ia diferencial: construye subvariedades
como preim\'agenes de valores regulares y es la base te\'orica de la estructura de
las fibras de una submersi\'on.

\begin{theorem}[Teorema del nivel regular]
\label{thm:app:regular-level}
Sea $F \colon \mathcal{M} \to \mathcal{N}$ una aplicaci\'on suave entre variedades
diferenciables de dimensiones $m = \dim \mathcal{M}$ y $n = \dim \mathcal{N}$, con
$m \geq n$. Si $q \in \mathcal{N}$ es un valor regular de $F$
(\Cref{def:app:regular-value}) con $F^{-1}(q) \neq \emptyset$, entonces
$F^{-1}(q)$ es una subvariedad suave de $\mathcal{M}$ de dimensi\'on $m - n$.
\end{theorem}

\begin{proof}
Sea $p \in F^{-1}(q)$. Como $q$ es un valor regular, el diferencial
$dF_p \colon T_p\mathcal{M} \to T_q\mathcal{N}$ es sobreyectivo, luego de rango
$n$. Elegimos cartas $(U, \varphi)$ centrada en $p$ y $(V, \psi)$ centrada en $q$
con $F(U) \subseteq V$. En estas coordenadas, la representaci\'on local
$\hat{F} = \psi \circ F \circ \varphi^{-1} \colon \varphi(U) \subseteq \R^m \to
\R^n$ tiene jacobiana $D\hat{F}_{\varphi(p)}$ de rango $n$.

Reordenando coordenadas si es necesario, podemos suponer que las \'ultimas $n$
columnas de $D\hat{F}_{\varphi(p)}$ forman un bloque $n \times n$ invertible.
Escribimos $\R^m = \R^{m-n} \times \R^n$ con coordenadas $(x', x'')$, donde
$x' = (x_1, \ldots, x_{m-n})$ y $x'' = (x_{m-n+1}, \ldots, x_m)$.

Consideramos la aplicaci\'on $\Phi \colon \varphi(U) \to \R^m$ dada por
$\Phi(x', x'') = (x', \hat{F}(x', x''))$. Su jacobiana en $\varphi(p)$ es
\[
D\Phi_{\varphi(p)} = \begin{pmatrix} I_{m-n} & 0 \\[2pt]
\frac{\partial \hat{F}}{\partial x'} & \frac{\partial \hat{F}}{\partial x''}
\end{pmatrix} ,
\]
que es invertible por ser triangular inferior por bloques con bloques diagonales
invertibles $I_{m-n}$ y $\frac{\partial \hat{F}}{\partial x''}$. Por el teorema de
la funci\'on inversa (\Cref{thm:app:inverse-function}), $\Phi$ es un difeomorfismo
local en un entorno $W$ de $\varphi(p)$.

En las nuevas coordenadas $\Phi$, la aplicaci\'on $F$ se lee como
\[
\hat{F} \circ \Phi^{-1}(y', y'') = y'' ,
\]
es decir, $\hat{F}$ se convierte en la proyecci\'on sobre las \'ultimas $n$
coordenadas. Por tanto,
\[
(\hat{F} \circ \Phi^{-1})^{-1}(\psi(q))
= \{(y', y'') \in \Phi(W) : y'' = \psi(q)\}
= \{(y', \psi(q)) : y' \in U'\}
\]
para cierto abierto $U' \subseteq \R^{m-n}$. Esto muestra que, cerca de cada punto
$p$, $F^{-1}(q)$ es localmente el grafo de una funci\'on suave de $m - n$
variables, lo que le confiere estructura de subvariedad suave de dimensi\'on
$m - n$.
\end{proof}

\begin{corollary}[Fibras de una submersi\'on]
\label{cor:app:submersion-fibres}
Si $F \colon \mathcal{M} \to \mathcal{N}$ es una submersi\'on suave, entonces para
cada $q \in F(\mathcal{M})$ la fibra $F^{-1}(q)$ es una subvariedad suave de
$\mathcal{M}$ de dimensi\'on $\dim \mathcal{M} - \dim \mathcal{N}$.
\end{corollary}

\begin{proof}
Como $F$ es una submersi\'on, $dF_p$ es sobreyectivo para todo
$p \in \mathcal{M}$. En particular, todo punto $q \in F(\mathcal{M})$ es un valor
regular de $F$, y el resultado se sigue del \Cref{thm:app:regular-level}.
\end{proof}

\begin{remark}
Una referencia completa para el teorema del nivel regular y sus aplicaciones puede
encontrarse en \citet[Cap\'itulo~5, Teorema~5.12]{lee2012introduction}. La
formulaci\'on cl\'asica en $\R^n$ se obtiene como caso particular tomando
$\mathcal{M} = \R^m$ y $\mathcal{N} = \R^n$ con las cartas identidad.
\end{remark}

%=======================================================================
\section{\'Area de Discos Geod\'esicos en Superficies de Curvatura Constante}
\label{sec:app:geodesic-disc}
%=======================================================================

El \'area de un disco geod\'esico de radio $R$ en una superficie de curvatura
gaussiana constante $K$ satisface una f\'ormula que generaliza el resultado
euclidiano $\pi R^2$. Derivamos esta f\'ormula a partir de la ecuaci\'on de Jacobi
en coordenadas polares geod\'esicas.

\begin{proposition}[M\'etrica en coordenadas polares geod\'esicas]
\label{prop:app:geodesic-polar-metric}
Sea $(\mathcal{M}, g)$ una superficie riemanniana de curvatura gaussiana constante
$K$, y sea $p \in \mathcal{M}$. En coordenadas polares geod\'esicas $(r, \theta)$
centradas en $p$ (donde $r = d_g(p, \cdot)$ y $\theta \in [0, 2\pi)$ es el
\'angulo en $T_p\mathcal{M}$), la m\'etrica riemanniana toma la forma
\begin{equation}
\label{eq:app:geodesic-polar-metric}
ds^2 = dr^2 + f_K(r)^2 \, d\theta^2 ,
\end{equation}
donde la funci\'on $f_K \colon [0, \infty) \to \R$ satisface la \textbf{ecuaci\'on
de Jacobi}
\begin{equation}
\label{eq:app:jacobi-radial}
f_K''(r) + K \, f_K(r) = 0 , \qquad f_K(0) = 0 , \quad f_K'(0) = 1 .
\end{equation}
\end{proposition}

\begin{proof}
En coordenadas polares geod\'esicas, las geod\'esicas radiales desde $p$
corresponden a $\theta = \text{const}$, con parametrizaci\'on natural por longitud
de arco $r$. El coeficiente m\'etrico $g_{rr} = 1$ refleja esta parametrizaci\'on.
El coeficiente angular $g_{\theta\theta} = f_K(r)^2$ mide c\'omo se separan las
geod\'esicas radiales adyacentes; la funci\'on $f_K(r)$ es la norma del campo de
Jacobi $J(r)$ a lo largo de una geod\'esica radial con condiciones iniciales
$J(0) = 0$ y $\norm{J'(0)} = 1$. La ecuaci\'on de Jacobi
$J'' + R(\dot{\gamma}, J)\dot{\gamma} = 0$, en el caso de curvatura constante $K$
en dimensi\'on $2$, se reduce a la ecuaci\'on escalar $f_K'' + K f_K = 0$
\citep[Cap\'itulo~10]{lee2018introduction}.
\end{proof}

La soluci\'on de la ecuaci\'on~\eqref{eq:app:jacobi-radial} depende del signo de
$K$:
\begin{equation}
\label{eq:app:fK-solutions}
f_K(r) = \begin{cases}
\dfrac{1}{\sqrt{K}} \sin(\sqrt{K} \, r) & \text{si } K > 0, \\[8pt]
r & \text{si } K = 0, \\[8pt]
\dfrac{1}{\sqrt{\abs{K}}} \sinh(\sqrt{\abs{K}} \, r) & \text{si } K < 0.
\end{cases}
\end{equation}
La verificaci\'on es inmediata: para $K > 0$,
$f_K(r) = \frac{1}{\sqrt{K}} \sin(\sqrt{K} r)$ satisface $f_K(0) = 0$,
$f_K'(0) = \cos(0) = 1$ y $f_K''(r) = -K f_K(r)$. Los casos $K = 0$ y $K < 0$ se
comprueban an\'alogamente.

\begin{theorem}[\'Area de un disco geod\'esico]
\label{thm:app:geodesic-disc-area}
Sea $(\mathcal{M}, g)$ una superficie de curvatura gaussiana constante $K$, y sea
$B_R(p) = \{q \in \mathcal{M} : d_g(p,q) \leq R\}$ el disco geod\'esico de radio
$R$ centrado en $p$ (con $R$ menor que el radio de inyectividad). Entonces
\begin{equation}
\label{eq:app:geodesic-disc-area}
A(R) = \begin{cases}
\dfrac{2\pi}{K}\bigl(1 - \cos(\sqrt{K} \, R)\bigr) & \text{si } K > 0, \\[8pt]
\pi R^2 & \text{si } K = 0, \\[8pt]
\dfrac{2\pi}{\abs{K}}\bigl(\cosh(\sqrt{\abs{K}} \, R) - 1\bigr) & \text{si } K < 0.
\end{cases}
\end{equation}
Adem\'as, $\lim_{K \to 0} A(R) = \pi R^2$ en todos los casos.
\end{theorem}

\begin{proof}
El elemento de \'area en coordenadas polares geod\'esicas es
$dA = f_K(r) \, dr \, d\theta$ (por~\eqref{eq:app:geodesic-polar-metric}, ya que
$\sqrt{\det(g)} = f_K(r)$). Por tanto,
\[
A(R) = \int_0^{2\pi} \int_0^R f_K(r) \, dr \, d\theta
= 2\pi \int_0^R f_K(r) \, dr .
\]

\textbf{Caso $K > 0$:}
\[
A(R) = 2\pi \int_0^R \frac{1}{\sqrt{K}} \sin(\sqrt{K} \, r) \, dr
= 2\pi \left[ \frac{-\cos(\sqrt{K} \, r)}{K} \right]_0^R
= \frac{2\pi}{K} \bigl(1 - \cos(\sqrt{K} \, R)\bigr) .
\]

\textbf{Caso $K = 0$:}
\[
A(R) = 2\pi \int_0^R r \, dr = \pi R^2 .
\]

\textbf{Caso $K < 0$:} escribiendo $K = -\abs{K}$,
\[
A(R) = 2\pi \int_0^R \frac{1}{\sqrt{\abs{K}}} \sinh(\sqrt{\abs{K}} \, r) \, dr
= 2\pi \left[ \frac{\cosh(\sqrt{\abs{K}} \, r)}{\abs{K}} \right]_0^R
= \frac{2\pi}{\abs{K}} \bigl(\cosh(\sqrt{\abs{K}} \, R) - 1\bigr) .
\]

\textbf{L\'imite $K \to 0$:} usando el desarrollo de Taylor
$\cos(x) = 1 - \frac{x^2}{2} + \frac{x^4}{24} - \cdots$,
\[
\frac{1 - \cos(\sqrt{K} R)}{K}
= \frac{1 - \bigl(1 - \frac{KR^2}{2} + \frac{K^2 R^4}{24} - \cdots\bigr)}{K}
= \frac{R^2}{2} - \frac{K R^4}{24} + \cdots
\xrightarrow{K \to 0} \frac{R^2}{2} .
\]
Multiplicando por $2\pi$ se obtiene $\pi R^2$. El caso $K < 0$ es an\'alogo usando
$\cosh(x) = 1 + \frac{x^2}{2} + \cdots$.
\end{proof}

\begin{remark}[Interpretaci\'on geom\'etrica]
\label{rem:app:geodesic-disc-interpretation}
El \Cref{thm:app:geodesic-disc-area} cuantifica el efecto de la curvatura sobre el
crecimiento del \'area:
\begin{itemize}
    \item \textbf{Curvatura positiva} ($K > 0$): $A(R) < \pi R^2$. El \'area crece
          m\'as lentamente que en el caso euclidiano. En la esfera $S^2$
          ($K = 1$), $A(R) = 2\pi(1 - \cos R)$; para $R = \pi/2$ (un cuarto de
          esfera), $A = 2\pi$.
    \item \textbf{Curvatura negativa} ($K < 0$): $A(R) > \pi R^2$. Como
          $\cosh(x) \sim \frac{1}{2}e^x$ para $x$ grande, el \'area crece
          \emph{exponencialmente}:
          $A(R) \sim \frac{\pi}{\abs{K}} e^{\sqrt{\abs{K}} R}$.
    \item \textbf{Conexi\'on con la ecuaci\'on de Jacobi.} La funci\'on $f_K(r)$
          que aparece en la m\'etrica es exactamente la que describe la
          desviaci\'on de geod\'esicas paralelas, mostrando que el crecimiento del
          \'area y la divergencia/convergencia de geod\'esicas son manifestaciones
          del mismo fen\'omeno geom\'etrico.
\end{itemize}
Una demostraci\'on alternativa puede encontrarse en
\citet[Cap\'itulo~4, Secci\'on~4.6]{do1976differential} para superficies y en
\citet[Teorema~11.9]{lee2018introduction} para el caso riemanniano general.
\end{remark}

%=======================================================================
\section{Complejidad Muestral en Variedades}
\label{sec:app:sample-complexity}
%=======================================================================

El n\'umero de muestras necesarias para aproximar una funci\'on Lipschitz sobre
una variedad escala con la dimensi\'on intr\'inseca $m$, no con la dimensi\'on
ambiente $d$. Presentamos los preliminares necesarios y un esbozo de la
demostraci\'on basado en los resultados de \citet{niyogi2008finding}.

\subsection{Preliminares}
\label{subsec:app:sample-complexity-prelim}

\begin{definition}[N\'umero de recubrimiento]
\label{def:app:covering-number}
Sea $(X, d)$ un espacio m\'etrico y $\epsilon > 0$. Un \emph{$\epsilon$-recubrimiento}
(o $\epsilon$-red) de $X$ es un subconjunto $S \subseteq X$ tal que para todo
$x \in X$ existe $s \in S$ con $d(x, s) \leq \epsilon$. El \emph{n\'umero de
recubrimiento} $\mathcal{N}(X, d, \epsilon)$ es el m\'inimo cardinal de un
$\epsilon$-recubrimiento de $X$.
\end{definition}

\begin{remark}
\label{rem:app:covering-vs-topological}
La noci\'on de $\epsilon$-recubrimiento m\'etrico no debe confundirse con la de
recubrimiento topol\'ogico. Un $\epsilon$-recubrimiento es un conjunto de
\emph{puntos} cuyas bolas de radio $\epsilon$ cubren el espacio, mientras que un
recubrimiento topol\'ogico es una familia de conjuntos abiertos cuya uni\'on
contiene al espacio. Ambos est\'an relacionados: si $S$ es un
$\epsilon$-recubrimiento de $(X,d)$, entonces $\{B(s, \epsilon)\}_{s \in S}$ es un
recubrimiento abierto de $X$ en sentido topol\'ogico.
\end{remark}

\begin{lemma}[Cota volum\'etrica del n\'umero de recubrimiento]
\label{lem:app:volumetric-covering}
Sea $(\mathcal{M}, g)$ una variedad riemanniana compacta y $\epsilon > 0$.
Entonces
\[
\mathcal{N}(\mathcal{M}, d_g, \epsilon) \geq
\frac{\Vol(\mathcal{M})}{\sup_{p \in \mathcal{M}} \Vol(B_g(p, \epsilon))} ,
\]
donde $B_g(p, \epsilon)$ denota la bola geod\'esica de centro $p$ y radio
$\epsilon$.
\end{lemma}

\begin{proof}
Sea $S = \{s_1, \ldots, s_N\}$ un $\epsilon$-recubrimiento \'optimo, es decir,
$\abs{S} = \mathcal{N}(\mathcal{M}, d_g, \epsilon)$. Por definici\'on,
$\mathcal{M} \subseteq \bigcup_{i=1}^N B_g(s_i, \epsilon)$, de modo que
\[
\Vol(\mathcal{M}) \leq \sum_{i=1}^N \Vol(B_g(s_i, \epsilon))
\leq N \cdot \sup_{p \in \mathcal{M}} \Vol(B_g(p, \epsilon)) .
\]
Despejando $N$ se obtiene la cota deseada.
\end{proof}

\begin{theorem}[Comparaci\'on de Bishop--G\"unther]
\label{thm:app:bishop-gunther}
Sea $(\mathcal{M}^m, g)$ una variedad riemanniana completa de dimensi\'on $m$ cuya
curvatura de Ricci satisface $\Ric \geq (m-1)\kappa_0 \, g$ para alg\'un
$\kappa_0 \in \R$. Sea $\mathbb{M}^m_{\kappa_0}$ el espacio modelo de curvatura
seccional constante $\kappa_0$. Entonces, para todo $p \in \mathcal{M}$ y $r > 0$
(con $r < \pi/\sqrt{\kappa_0}$ si $\kappa_0 > 0$),
\[
\Vol(B_g(p, r)) \leq \Vol(B_{\kappa_0}(r)) ,
\]
donde $B_{\kappa_0}(r)$ denota la bola geod\'esica de radio $r$ en
$\mathbb{M}^m_{\kappa_0}$.
\end{theorem}

\begin{remark}
\label{rem:app:bishop-gunther-context}
Este teorema generaliza a dimensi\'on arbitraria la comparaci\'on de \'areas del
\Cref{thm:app:geodesic-disc-area}, que trata el caso bidimensional. En el espacio
modelo, el volumen de la bola de radio $r$ es
$\Vol(B_{\kappa_0}(r)) = \omega_m \int_0^r \operatorname{sn}_{\kappa_0}(t)^{m-1}
\, dt$, donde $\omega_m$ es el volumen de la bola unitaria en $\R^m$ y
$\operatorname{sn}_{\kappa_0}$ es la funci\'on de comparaci\'on
($\sin(\sqrt{\kappa_0}\,t)/\sqrt{\kappa_0}$ si $\kappa_0 > 0$, $t$ si
$\kappa_0 = 0$, $\sinh(\sqrt{\abs{\kappa_0}}\,t)/\sqrt{\abs{\kappa_0}}$ si
$\kappa_0 < 0$). Para radios peque\~nos,
$\Vol(B_{\kappa_0}(r)) = \omega_m r^m (1 + O(\kappa_0 r^2))$. Para una
demostraci\'on completa, v\'ease \citet[Teorema~11.10]{lee2018introduction}.
\end{remark}

\subsection{Esbozo de la demostraci\'on}
\label{subsec:app:sample-complexity-proof}

Esbozamos la demostraci\'on de la cota de complejidad muestral en variedades en
tres pasos.

\paragraph{Paso 1: Reducci\'on a n\'umeros de recubrimiento.}
Sea $f \colon \mathcal{M} \to \R$ una funci\'on $L$-Lipschitz y sea
$\{x_1, \ldots, x_N\} \subset \mathcal{M}$ un conjunto de muestras. Si queremos
aproximar $f$ uniformemente con error $\epsilon$ usando solo los valores
$f(x_1), \ldots, f(x_N)$, entonces para cualquier punto $q \in \mathcal{M}$ debe
existir una muestra $x_i$ con $d_g(q, x_i) \leq \epsilon/L$. En efecto, si
$d_g(q, x_i) > \epsilon/L$ para todo $i$, la condici\'on Lipschitz solo garantiza
$\abs{f(q) - f(x_i)} \leq L \cdot d_g(q, x_i)$, que no puede acotarse por
$\epsilon$. Por tanto, las muestras deben formar un $(\epsilon/L)$-recubrimiento de
$\mathcal{M}$, y se requieren al menos
$\mathcal{N}(\mathcal{M}, d_g, \epsilon/L)$ muestras.

\paragraph{Paso 2: Cota inferior del n\'umero de recubrimiento.}
Aplicando el \Cref{lem:app:volumetric-covering} con radio $\epsilon/L$,
\[
N \geq \mathcal{N}\!\left(\mathcal{M}, d_g, \frac{\epsilon}{L}\right)
\geq \frac{\Vol(\mathcal{M})}{\sup_{p} \Vol\!\left(B_g\!\left(p,
\frac{\epsilon}{L}\right)\right)} .
\]

\paragraph{Paso 3: Estimaci\'on de vol\'umenes.}
La curvatura seccional de $\mathcal{M}$ est\'a acotada en valor absoluto por
$\kappa$, lo que implica $\Ric \geq -(m-1)\kappa \, g$. Por el teorema de
Bishop--G\"unther (\Cref{thm:app:bishop-gunther}) con $\kappa_0 = -\kappa$,
\[
\Vol\!\left(B_g\!\left(p, \frac{\epsilon}{L}\right)\right)
\leq \omega_m \left(\frac{\epsilon}{L}\right)^m
\left(1 + O\!\left(\kappa \cdot \frac{\epsilon^2}{L^2}\right)\right) .
\]
El volumen total $\Vol(\mathcal{M})$ admite una cota inferior que depende del
alcance (\emph{reach}) $\tau$ y de $\kappa$ (por resultados de
\citet{niyogi2008finding} sobre la geometr\'ia del \emph{reach}). Combinando,
\[
N \geq \frac{\Vol(\mathcal{M})}{\omega_m (\epsilon/L)^m
(1 + O(\kappa \epsilon^2/L^2))}
= \Omega\!\left(\left(\frac{L}{\epsilon}\right)^m \cdot
\operatorname{poly}(\tau^{-1}, \kappa)\right) .
\]
Absorbiendo la constante de Lipschitz $L$ en el factor $\operatorname{poly}$ se
obtiene la cota.

\begin{remark}[Dimensi\'on intr\'inseca frente a ambiente]
\label{rem:app:intrinsic-complexity}
El resultado clave es que el exponente en $1/\epsilon$ es la dimensi\'on
intr\'inseca $m$, no la dimensi\'on ambiente $d$. Para datos en $\R^d$ que residen
en una variedad $m$-dimensional con $m \ll d$, la cota $\Omega((1/\epsilon)^m)$ es
exponencialmente menor que la cota $\Omega((1/\epsilon)^d)$ que se obtendr\'ia
ignorando la estructura geom\'etrica. Esta es la formalizaci\'on matem\'atica de
la maldici\'on de la dimensionalidad atenuada por la hip\'otesis de la variedad.
\end{remark}

%=======================================================================
\section{Isometr\'ias entre Modelos del Espacio Hiperb\'olico}
\label{sec:app:hyperbolic-isometries}
%=======================================================================

Existen tres modelos del espacio hiperb\'olico (hiperboloide, disco de Poincar\'e,
semiplano superior), y son isom\'etricos como variedades riemannianas. Demostramos
expl\'icitamente las isometr\'ias entre ellos mediante la proyecci\'on
estereogr\'afica y la transformada de Cayley.

\subsection{Preliminares}
\label{subsec:app:hyperbolic-prelim}

\begin{definition}[Isometr\'ia riemanniana]
\label{def:app:riemannian-isometry}
Sean $(\mathcal{M}, g)$ y $(\mathcal{N}, h)$ variedades riemannianas. Un
difeomorfismo $F \colon \mathcal{M} \to \mathcal{N}$ es una \emph{isometr\'ia
riemanniana} si preserva la m\'etrica, es decir, si $F^*h = g$, donde el
\emph{pullback} de la m\'etrica se define por
\begin{equation}
\label{eq:app:metric-pullback}
(F^*h)_p(v, w) = h_{F(p)}(dF_p(v), dF_p(w))
\end{equation}
para todo $p \in \mathcal{M}$ y $v, w \in T_p\mathcal{M}$.
\end{definition}

\begin{remark}
\label{rem:app:conformal-isometry}
Cuando ambas m\'etricas son \emph{conformes}, es decir, de la forma
$g_{ij} = \lambda(x)^2 \delta_{ij}$ y $h_{ij} = \mu(y)^2 \delta_{ij}$, verificar la
condici\'on $F^*h = g$ se simplifica: basta comprobar que
\[
\mu(F(x))^2 \cdot \norm{dF_x(v)}^2 = \lambda(x)^2 \cdot \norm{v}^2
\]
para todo vector tangente $v$, donde $\norm{\cdot}$ es la norma euclidiana. Esto
reduce la verificaci\'on a una identidad escalar entre los factores conformes.
V\'ease \citet[Cap\'itulo~3]{lee2018introduction}.
\end{remark}

\subsection{Del hiperboloide al disco de Poincar\'e}
\label{subsec:app:hyperboloid-poincare}

\begin{theorem}[Isometr\'ia hiperboloide--disco de Poincar\'e]
\label{thm:app:hyperboloid-poincare}
La proyecci\'on estereogr\'afica desde el polo sur del hiperboloide,
\begin{equation}
\label{eq:app:stereographic-hyperboloid}
\pi \colon \mathbb{H}^n \to \mathbb{D}^n , \qquad
\pi(x_0, x_1, \ldots, x_n) = \frac{(x_1, \ldots, x_n)}{1 + x_0} ,
\end{equation}
es una isometr\'ia riemanniana
$(\mathbb{H}^n, g_{\mathcal{L}}) \to (\mathbb{D}^n, g_{\mathbb{D}})$, donde
$g_{\mathcal{L}}$ es la m\'etrica inducida por el producto de Lorentz sobre el
hiperboloide y $g_{\mathbb{D}} = \frac{4}{(1-\norm{\mathbf{y}}^2)^2}\delta_{ij}$ es
la m\'etrica de Poincar\'e. Su inversa es
\begin{equation}
\label{eq:app:inverse-stereographic}
\pi^{-1}(y_1, \ldots, y_n) = \frac{1}{1 - \norm{\mathbf{y}}^2}
\bigl(1 + \norm{\mathbf{y}}^2,\; 2y_1,\; \ldots,\; 2y_n\bigr) .
\end{equation}
\end{theorem}

\begin{proof}
Procedemos en varios pasos.

\textbf{Paso 1: Interpretaci\'on geom\'etrica.}
La aplicaci\'on $\pi$ es la proyecci\'on estereogr\'afica desde el punto
$(-1, 0, \ldots, 0)$ (el polo sur del hiperboloide en el espacio de Minkowski)
sobre el hiperplano $\{x_0 = 0\} \cong \R^n$. La recta que une $(-1, \mathbf{0})$
con $(x_0, \mathbf{x}) \in \mathbb{H}^n$ es
$t \mapsto (-1, \mathbf{0}) + t\bigl((x_0+1), \mathbf{x}\bigr)$, y corta al
hiperplano $x_0 = 0$ cuando $t = \frac{1}{1+x_0}$, dando el punto
$\frac{\mathbf{x}}{1+x_0}$.

\textbf{Paso 2: Biyectividad y suavidad.}
Tanto $\pi$ como $\pi^{-1}$ son funciones racionales de sus argumentos, luego
suaves en sus dominios. Verificamos que $\pi^{-1}$ est\'a bien definida: si
$\norm{\mathbf{y}} < 1$, debemos comprobar $\pi^{-1}(\mathbf{y}) \in \mathbb{H}^n$.
Sea $s = \norm{\mathbf{y}}^2 < 1$. Entonces
\[
\inner{\pi^{-1}(\mathbf{y})}{\pi^{-1}(\mathbf{y})}_{\mathcal{L}}
= \frac{1}{(1-s)^2}\bigl(-(1+s)^2 + 4s\bigr)
= \frac{-(1-s)^2}{(1-s)^2} = -1 ,
\]
y la componente $x_0 = \frac{1+s}{1-s} > 0$, de modo que
$\pi^{-1}(\mathbf{y}) \in \mathbb{H}^n$.

Para la composici\'on $\pi \circ \pi^{-1}$: dado $\mathbf{y} \in \mathbb{D}^n$,
$\pi^{-1}(\mathbf{y})$ tiene componente
$x_0 = \frac{1+\norm{\mathbf{y}}^2}{1-\norm{\mathbf{y}}^2}$ y componentes
espaciales $\frac{2y_i}{1-\norm{\mathbf{y}}^2}$. Entonces
\[
\pi(\pi^{-1}(\mathbf{y}))
= \frac{(2y_1, \ldots, 2y_n)/(1-\norm{\mathbf{y}}^2)}
{1 + (1+\norm{\mathbf{y}}^2)/(1-\norm{\mathbf{y}}^2)}
= \frac{(2y_1, \ldots, 2y_n)/(1-\norm{\mathbf{y}}^2)}
{2/(1-\norm{\mathbf{y}}^2)} = \mathbf{y} .
\]
La verificaci\'on de $\pi^{-1} \circ \pi = \mathrm{id}$ es an\'aloga.

\textbf{Paso 3: Relaci\'on entre norma y coordenadas.}
Sea $(x_0, \mathbf{x}) \in \mathbb{H}^n$ y
$\mathbf{y} = \pi(x_0, \mathbf{x}) = \frac{\mathbf{x}}{1+x_0}$. Usando la relaci\'on
del hiperboloide $x_0^2 - \norm{\mathbf{x}}^2 = 1$,
\begin{equation}
\label{eq:app:norm-y-hyperboloid}
\norm{\mathbf{y}}^2 = \frac{\norm{\mathbf{x}}^2}{(1+x_0)^2}
= \frac{x_0^2 - 1}{(1+x_0)^2} = \frac{x_0 - 1}{x_0 + 1} .
\end{equation}
De aqu\'i se deduce la identidad clave
\begin{equation}
\label{eq:app:conformal-factor-key}
1 - \norm{\mathbf{y}}^2 = 1 - \frac{x_0 - 1}{x_0 + 1} = \frac{2}{1 + x_0} .
\end{equation}

\textbf{Paso 4: C\'alculo del diferencial.}
Para calcular $d\pi$, usamos coordenadas intr\'insecas $(x_1, \ldots, x_n)$ en
$\mathbb{H}^n$. La componente $x_0$ satisface
$x_0 = \sqrt{1 + \norm{\mathbf{x}}^2}$, de donde
\[
dx_0 = \frac{\sum_{i=1}^n x_i \, dx_i}{x_0} .
\]
Diferenciando $y_j = \frac{x_j}{1+x_0}$,
\begin{align}
dy_j &= \frac{dx_j(1+x_0) - x_j \, dx_0}{(1+x_0)^2} \nonumber \\
&= \frac{dx_j}{1+x_0} - \frac{x_j}{(1+x_0)^2} \cdot
\frac{\sum_i x_i \, dx_i}{x_0} . \label{eq:app:differential-pi}
\end{align}

\textbf{Paso 5: Verificaci\'on de la isometr\'ia.}
Debemos comprobar que $\pi^* g_{\mathbb{D}} = g_{\mathcal{L}}$. La m\'etrica
inducida por Lorentz en $\mathbb{H}^n$ con coordenadas intr\'insecas
$(x_1, \ldots, x_n)$ es
\[
g_{\mathcal{L}} = -dx_0^2 + \sum_{i=1}^n dx_i^2
= -\frac{(\sum_i x_i \, dx_i)^2}{x_0^2} + \sum_{i=1}^n dx_i^2 ,
\]
es decir, en componentes
$(g_{\mathcal{L}})_{ij} = \delta_{ij} - \frac{x_i x_j}{x_0^2}$. El pullback de la
m\'etrica de Poincar\'e es
\[
(\pi^* g_{\mathbb{D}})_{ij} = \frac{4}{(1-\norm{\mathbf{y}}^2)^2}
\sum_{k=1}^n \frac{\partial y_k}{\partial x_i} \frac{\partial y_k}{\partial x_j} .
\]
Usando~\eqref{eq:app:conformal-factor-key}, el factor conforme es
$\frac{4}{(1-\norm{\mathbf{y}}^2)^2} = (1+x_0)^2$.
De~\eqref{eq:app:differential-pi},
\[
\frac{\partial y_k}{\partial x_i} = \frac{\delta_{ki}}{1+x_0}
- \frac{x_k x_i}{x_0(1+x_0)^2} ,
\]
de modo que
\begin{align*}
\sum_{k} \frac{\partial y_k}{\partial x_i}\frac{\partial y_k}{\partial x_j}
&= \frac{\delta_{ij}}{(1+x_0)^2}
- \frac{2 x_i x_j}{x_0(1+x_0)^3}
+ \frac{x_i x_j (x_0^2 - 1)}{x_0^2(1+x_0)^4} \\
&= \frac{\delta_{ij}}{(1+x_0)^2}
- \frac{2 x_i x_j}{x_0(1+x_0)^3}
+ \frac{x_i x_j (x_0 - 1)}{x_0^2(1+x_0)^3} \\
&= \frac{\delta_{ij}}{(1+x_0)^2} - \frac{x_i x_j}{x_0^2(1+x_0)^2} ,
\end{align*}
donde en el \'ultimo paso se utiliza
\[
-\frac{2}{x_0(1+x_0)} + \frac{x_0-1}{x_0^2(1+x_0)}
= \frac{-2x_0 + x_0 - 1}{x_0^2(1+x_0)}
= \frac{-(x_0+1)}{x_0^2(1+x_0)} = \frac{-1}{x_0^2} .
\]
Multiplicando por $(1+x_0)^2$,
\[
(\pi^* g_{\mathbb{D}})_{ij} = (1+x_0)^2 \left(\frac{\delta_{ij}}{(1+x_0)^2}
- \frac{x_i x_j}{x_0^2(1+x_0)^2}\right)
= \delta_{ij} - \frac{x_i x_j}{x_0^2} = (g_{\mathcal{L}})_{ij} .
\]
Esto demuestra que $\pi^* g_{\mathbb{D}} = g_{\mathcal{L}}$, es decir, $\pi$ es una
isometr\'ia.
\end{proof}

\subsection{Del disco de Poincar\'e al semiplano superior}
\label{subsec:app:poincare-halfplane}

\begin{theorem}[Isometr\'ia disco de Poincar\'e--semiplano superior]
\label{thm:app:poincare-halfplane}
En dimensi\'on $n = 2$, la transformada de Cayley
\begin{equation}
\label{eq:app:cayley-transform}
C \colon \mathbb{D}^2 \to \mathbb{H}^2 , \qquad C(w) = i\,\frac{1+w}{1-w} ,
\end{equation}
es una isometr\'ia riemanniana del disco de Poincar\'e
$(\mathbb{D}^2, g_{\mathbb{D}})$ al semiplano superior
$(\mathbb{H}^2, g_{\mathbb{H}})$, donde
$g_{\mathbb{H}} = \frac{dx^2 + dy^2}{y^2}$. Su inversa es
\begin{equation}
\label{eq:app:cayley-inverse}
C^{-1} \colon \mathbb{H}^2 \to \mathbb{D}^2 , \qquad C^{-1}(z) = \frac{z - i}{z + i} .
\end{equation}
\end{theorem}

\begin{proof}
Procedemos en cuatro pasos.

\textbf{Paso 1: Buena definici\'on ($C(\mathbb{D}^2) \subseteq \mathbb{H}^2$).}
Sea $w \in \mathbb{D}^2$, es decir, $\abs{w} < 1$. Debemos verificar
$\Img(C(w)) > 0$. Un c\'alculo directo da
\[
\frac{1+w}{1-w} = \frac{(1+w)(1-\bar{w})}{\abs{1-w}^2}
= \frac{1 - \bar{w} + w - \abs{w}^2}{\abs{1-w}^2} ,
\]
de donde
\[
\Img(C(w)) = \operatorname{Re}\!\left(\frac{1+w}{1-w}\right)
= \frac{1 - \abs{w}^2}{\abs{1-w}^2} .
\]
Como $\abs{w} < 1$, tenemos $1 - \abs{w}^2 > 0$ y $\abs{1-w}^2 > 0$, por lo que
$\Img(C(w)) > 0$.

\textbf{Paso 2: Biyectividad.}
Para $C^{-1}(z) = \frac{z-i}{z+i}$ con $\Img(z) > 0$, escribiendo $z = a + bi$ con
$b > 0$,
\[
\abs{C^{-1}(z)}^2 = \frac{\abs{z-i}^2}{\abs{z+i}^2}
= \frac{a^2 + (b-1)^2}{a^2 + (b+1)^2}
= \frac{a^2 + b^2 - 2b + 1}{a^2 + b^2 + 2b + 1} < 1 ,
\]
ya que $-2b < 2b$ cuando $b > 0$. Luego $C^{-1}(z) \in \mathbb{D}^2$. Las
composiciones $C \circ C^{-1} = \mathrm{id}$ y $C^{-1} \circ C = \mathrm{id}$ se
verifican por c\'alculo directo.

\textbf{Paso 3: Diferencial.}
La derivada compleja de $C$ es
\begin{equation}
\label{eq:app:cayley-derivative}
C'(w) = i \cdot \frac{(1-w) - (1+w)(-1)}{(1-w)^2} = \frac{2i}{(1-w)^2} .
\end{equation}

\textbf{Paso 4: Verificaci\'on de la isometr\'ia.}
Para una aplicaci\'on holomorfa $C$ entre superficies con m\'etricas conformes, el
pullback se calcula mediante
\[
(C^* g_{\mathbb{H}})_w = \frac{\abs{C'(w)}^2}{(\Img C(w))^2} \cdot (du^2 + dv^2) ,
\]
donde $w = u + iv$. Sustituyendo los resultados de los pasos anteriores,
\[
\frac{\abs{C'(w)}^2}{(\Img C(w))^2}
= \frac{4/\abs{1-w}^4}{(1-\abs{w}^2)^2/\abs{1-w}^4}
= \frac{4}{(1-\abs{w}^2)^2} ,
\]
que es exactamente el factor conforme de la m\'etrica de Poincar\'e
$g_{\mathbb{D}} = \frac{4}{(1-\abs{w}^2)^2}(du^2 + dv^2)$. Por tanto
$C^* g_{\mathbb{H}} = g_{\mathbb{D}}$.
\end{proof}

\subsection{Equivalencia de los tres modelos}
\label{subsec:app:three-models}

\begin{corollary}[Equivalencia de los tres modelos hiperb\'olicos]
\label{cor:app:hyperbolic-equivalence}
Los tres modelos del espacio hiperb\'olico (hiperboloide, disco de Poincar\'e y
semiplano superior) son riemannianamente isom\'etricos. En particular, la
composici\'on $C \circ \pi$ proporciona una isometr\'ia expl\'icita del modelo del
hiperboloide al semiplano superior.
\end{corollary}

\begin{proof}
Composici\'on de las isometr\'ias de los
\Cref{thm:app:hyperboloid-poincare,thm:app:poincare-halfplane}: como la
composici\'on de isometr\'ias es una isometr\'ia,
$C \circ \pi \colon (\mathbb{H}^n, g_{\mathcal{L}}) \to
(\mathbb{D}^n, g_{\mathbb{D}}) \to (\mathbb{H}^2, g_{\mathbb{H}})$ es una
isometr\'ia (en dimensi\'on $n = 2$ para el semiplano).
\end{proof}

\begin{remark}[Implicaciones para el aprendizaje profundo geom\'etrico]
\label{rem:app:models-gdl}
La equivalencia isom\'etrica de los tres modelos justifica tratar la elecci\'on
del modelo en aplicaciones de GDL como una cuesti\'on de conveniencia
computacional:
\begin{itemize}
    \item El \textbf{disco de Poincar\'e} $\mathbb{D}^n$ es preferido para la
          optimizaci\'on por gradiente, ya que los mapas exponencial y
          logar\'itmico admiten expresiones en forma cerrada~\citep{nickel2017poincare}.
    \item El \textbf{modelo del hiperboloide} ofrece mayor estabilidad
          num\'erica, puesto que el factor conforme
          $\frac{4}{(1-\norm{\mathbf{x}}^2)^2}$ del disco de Poincar\'e diverge
          cuando $\norm{\mathbf{x}} \to 1$ (cerca del borde del disco), mientras
          que la m\'etrica de Lorentz inducida no presenta esta singularidad.
    \item La transformaci\'on $\pi$ del \Cref{thm:app:hyperboloid-poincare} es
          precisamente el mapeo utilizado en la pr\'actica para convertir entre
          representaciones.
\end{itemize}
\end{remark}

%=======================================================================
\section{Curvatura Seccional de los Espacios Modelo}
\label{sec:app:model-curvature}
%=======================================================================

La esfera $S^n$ tiene curvatura seccional constante $+1$ y el espacio
hiperb\'olico $\mathbb{H}^n$ tiene curvatura seccional constante $-1$. Demostramos
ambos resultados aplicando la definici\'on de curvatura seccional a las m\'etricas
expl\'icitas de estos espacios.

\subsection{Preliminares}
\label{subsec:app:model-curvature-prelim}

\begin{lemma}[Ecuaci\'on de Gauss para hipersuperficies]
\label{lem:app:gauss-equation}
Sea $M^n \subset \R^{n+1}$ una hipersuperficie orientada con normal unitaria
$\nu$, y sea $S_p \colon T_pM \to T_pM$ el \emph{operador de forma} (o de
Weingarten) $S_p(v) = -d\nu_p(v)$. La segunda forma fundamental es
$\II(u, v) = \inner{S_p(u)}{v}$. Si $u, v \in T_pM$ son ortonormales, entonces la
curvatura seccional de $M$ en el plano generado por $\{u, v\}$ satisface
\begin{equation}
\label{eq:app:gauss-equation}
K_M(u,v) = K_{\R^{n+1}}(u,v) + \II(u,u)\,\II(v,v) - \II(u,v)^2 .
\end{equation}
Como $\R^{n+1}$ es plano ($K_{\R^{n+1}} \equiv 0$), se reduce a
\begin{equation}
\label{eq:app:gauss-simplified}
K_M(u,v) = \II(u,u)\,\II(v,v) - \II(u,v)^2 .
\end{equation}
\end{lemma}

La demostraci\'on puede encontrarse en \citet[Teorema~8.4]{lee2018introduction} o
en \citet[Cap\'itulo~3, Proposici\'on~4]{do1976differential}.

\begin{lemma}[Curvatura de m\'etricas conformes en dimensi\'on 2]
\label{lem:app:conformal-curvature-2d}
Sea $g = \lambda(x,y)^2(dx^2 + dy^2)$ una m\'etrica conforme en un dominio de
$\R^2$, donde $\lambda > 0$ es suave. Entonces la curvatura gaussiana de
$(\R^2, g)$ es
\begin{equation}
\label{eq:app:conformal-curvature}
K = -\frac{\Delta(\ln \lambda)}{\lambda^2} ,
\end{equation}
donde $\Delta = \partial_{xx} + \partial_{yy}$ es el laplaciano euclidiano.
\end{lemma}

\begin{proof}
Para la m\'etrica $g_{ij} = \lambda^2 \delta_{ij}$, los s\'imbolos de Christoffel
son
\[
\Gamma^k_{ij} = (\partial_i \ln\lambda)\,\delta_{jk}
+ (\partial_j \ln\lambda)\,\delta_{ik} - (\partial_k \ln\lambda)\,\delta_{ij} .
\]
Sustituyendo en la expresi\'on del tensor de Riemann $R^l{}_{ijk}$ y contrayendo
para obtener la curvatura gaussiana $K = R^1{}_{212}/\lambda^2$, un c\'alculo
directo da~\eqref{eq:app:conformal-curvature}. Los detalles pueden consultarse en
\citet[Proposici\'on~8.32]{lee2018introduction}.
\end{proof}

\subsection{Curvatura de la esfera \texorpdfstring{$S^n$}{Sn}}
\label{subsec:app:sphere-curvature}

\begin{theorem}[Curvatura seccional de la esfera]
\label{thm:app:sphere-curvature}
La esfera unitaria $S^n \subset \R^{n+1}$ con la m\'etrica riemanniana inducida
por la m\'etrica euclidiana ambiente tiene curvatura seccional constante
$K = +1$.
\end{theorem}

\begin{proof}
Consideramos $S^n = \{p \in \R^{n+1} : \norm{p} = 1\}$ con la m\'etrica inducida
por $\R^{n+1}$.

\textbf{Paso 1: Normal unitaria.}
El campo normal unitario exterior es $\nu(p) = p$ para todo $p \in S^n$.

\textbf{Paso 2: Operador de forma.}
Para $v \in T_pS^n$, consideramos la curva
$\gamma(t) = \frac{p + tv}{\norm{p + tv}}$ en $S^n$ con $\gamma(0) = p$ y
$\gamma'(0) = v$ (pues $\inner{p}{v} = 0$). Entonces
$\nu(\gamma(t)) = \gamma(t) = \frac{p + tv}{\norm{p + tv}}$, y derivando en
$t = 0$,
\[
d\nu_p(v) = \frac{d}{dt}\bigg|_{t=0} \frac{p + tv}{\norm{p + tv}} = v ,
\]
ya que $\norm{p + tv}^2 = 1 + t^2\norm{v}^2$ implica
$\frac{d}{dt}\norm{p+tv}\big|_{t=0} = 0$. Por tanto, el operador de forma es
$S_p = -d\nu_p = -\operatorname{Id}_{T_pS^n}$.

\textbf{Paso 3: Segunda forma fundamental.}
Para $u, v \in T_pS^n$ ortonormales,
\[
\II(u,u) = \inner{S_p(u)}{u} = \inner{-u}{u} = -1 , \quad
\II(v,v) = -1 , \quad \II(u,v) = \inner{-u}{v} = 0 .
\]

\textbf{Paso 4: Ecuaci\'on de Gauss.}
Aplicando el \Cref{lem:app:gauss-equation},
\[
K(u,v) = (-1)(-1) - 0^2 = 1 .
\]
Como la elecci\'on de $p \in S^n$ y del $2$-plano $\{u,v\}$ es arbitraria,
$K \equiv +1$.
\end{proof}

\subsection{Curvatura del espacio hiperb\'olico \texorpdfstring{$\mathbb{H}^n$}{Hn}}
\label{subsec:app:hyperbolic-curvature}

\begin{theorem}[Curvatura seccional del espacio hiperb\'olico]
\label{thm:app:hyperbolic-curvature}
El disco de Poincar\'e $(\mathbb{D}^n, g_{\mathbb{D}})$ con la m\'etrica
$g_{\mathbb{D}} = \frac{4}{(1-\norm{\mathbf{x}}^2)^2}\delta_{ij}$ tiene curvatura
seccional constante $K = -1$.
\end{theorem}

\begin{proof}
Procedemos en dos partes: primero un c\'alculo expl\'icito de la curvatura en
dimensi\'on $2$, luego una extensi\'on a dimensi\'on arbitraria.

\textbf{Parte 1: Dimensi\'on 2 (c\'alculo expl\'icito).}
En $\mathbb{D}^2$ la m\'etrica es conforme con factor
$\lambda(x,y) = \frac{2}{1 - r^2}$, donde $r^2 = x^2 + y^2$. Calculamos
$\ln\lambda = \ln 2 - \ln(1 - r^2)$ y
\[
\partial_x(\ln\lambda) = \frac{2x}{1 - r^2} , \qquad
\partial_{xx}(\ln\lambda) = \frac{2(1-r^2) + 4x^2}{(1-r^2)^2} .
\]
Por simetr\'ia, $\partial_{yy}(\ln\lambda) = \frac{2(1-r^2) + 4y^2}{(1-r^2)^2}$.
Sumando,
\[
\Delta(\ln\lambda) = \frac{4(1-r^2) + 4r^2}{(1-r^2)^2} = \frac{4}{(1-r^2)^2} .
\]
Aplicando el \Cref{lem:app:conformal-curvature-2d},
\[
K = -\frac{\Delta(\ln\lambda)}{\lambda^2}
= -\frac{4/(1-r^2)^2}{4/(1-r^2)^2} = -1 .
\]

\textbf{Parte 2: Extensi\'on a dimensi\'on $n$ (argumento de simetr\'ia).}
La m\'etrica $g_{\mathbb{D}}$ es invariante bajo la acci\'on de $\Orth(n)$ por
rotaciones de $\mathbf{x}$ (pues $g_{\mathbb{D}}$ depende solo de
$\norm{\mathbf{x}}$). En cualquier punto $p \in \mathbb{D}^n$, el estabilizador de
$p$ bajo esta acci\'on act\'ua transitivamente sobre los $2$-planos de
$T_p\mathbb{D}^n$ que contienen la direcci\'on radial, y por composici\'on con
rotaciones, sobre todos los $2$-planos. Esto implica que todas las curvaturas
seccionales en $p$ son iguales.

El subespacio $\Pi = \{x \in \mathbb{D}^n : x_3 = \cdots = x_n = 0\}$ es una
subvariedad totalmente geod\'esica isom\'etrica a $(\mathbb{D}^2, g_{\mathbb{D}^2})$:
la reflexi\'on $\sigma \colon x_i \mapsto -x_i$ para $i \geq 3$ es una isometr\'ia
de $(\mathbb{D}^n, g_{\mathbb{D}})$ cuyo conjunto de puntos fijos es $\Pi$, y el
conjunto de puntos fijos de una isometr\'ia es totalmente geod\'esico
\citep[Proposici\'on~5.21]{lee2018introduction}. Por la Parte~1, la curvatura
seccional en $\Pi$ es $-1$, luego $K = -1$ para todos los $2$-planos en cada punto
de $\Pi$.

Finalmente, la m\'etrica de Poincar\'e es homog\'enea: las transformaciones de
M\"obius proporcionan isometr\'ias que llevan cualquier punto de $\mathbb{D}^n$ a
cualquier otro \citep{nickel2017poincare}. Por tanto, $K = -1$ en todo
$\mathbb{D}^n$.
\end{proof}

\begin{corollary}
\label{cor:app:hyperbolic-models-curvature}
Por los \Cref{thm:app:hyperboloid-poincare,thm:app:poincare-halfplane}, el modelo
del hiperboloide $(\mathbb{H}^n, g_{\mathcal{L}})$ y el semiplano superior
$(\mathbb{H}^2, g_{\mathbb{H}})$ tambi\'en tienen curvatura seccional constante
$K = -1$.
\end{corollary}

\begin{proof}
Las isometr\'ias preservan la curvatura seccional: si
$F \colon (M, g) \to (N, h)$ es una isometr\'ia riemanniana, entonces
$K_M(\sigma) = K_N(dF(\sigma))$ para todo $2$-plano $\sigma \subset T_pM$. El
resultado se sigue del \Cref{thm:app:hyperbolic-curvature} y de las isometr\'ias
expl\'icitas de la \Cref{sec:app:hyperbolic-isometries}.
\end{proof}

%=======================================================================
\section{El Teorema de Lagrange}
\label{sec:app:lagrange}
%=======================================================================

El teorema de Lagrange es un resultado fundamental de la teor\'ia de grupos
finitos que relaciona el orden de un grupo con los \'ordenes de sus subgrupos.
Aunque su demostraci\'on es elemental, incluimos los detalles por completitud.

\begin{theorem}[Lagrange]
\label{thm:app:lagrange}
Sea $\group$ un grupo finito y $H \leq \group$ un subgrupo. Entonces $\abs{H}$
divide a $\abs{\group}$, y
\begin{equation}
\label{eq:app:lagrange}
\abs{\group} = \abs{H} \cdot [\group:H] ,
\end{equation}
donde $[\group:H]$ denota el \'indice de $H$ en $\group$ (el n\'umero de clases
laterales izquierdas de $H$ en $\group$).
\end{theorem}

\begin{proof}
Definimos una relaci\'on de equivalencia en $\group$ por $a \sim b$ si y solo si
$a^{-1}b \in H$. Las clases de equivalencia son las clases laterales izquierdas
$g H = \{gh : h \in H\}$ para $g \in \group$.

\textbf{Paso 1.} Cada clase lateral tiene exactamente $\abs{H}$ elementos. En
efecto, la aplicaci\'on $\varphi_g \colon H \to gH$ dada por $\varphi_g(h) = gh$
es biyectiva: sobreyectiva por construcci\'on, e inyectiva porque $gh_1 = gh_2$
implica $h_1 = h_2$ por la ley de cancelaci\'on.

\textbf{Paso 2.} Las clases laterales forman una partici\'on de $\group$. Es
consecuencia directa de que $\sim$ es una relaci\'on de equivalencia.

\textbf{Paso 3.} Si $[\group:H]$ denota el n\'umero de clases laterales distintas,
entonces
\[
\abs{\group} = [\group:H] \cdot \abs{H} ,
\]
pues $\group$ se descompone en exactamente $[\group:H]$ conjuntos disjuntos, cada
uno de cardinal $\abs{H}$. En particular, $\abs{H}$ divide a $\abs{\group}$.
\end{proof}

%=======================================================================
\section{El Teorema de Rad\'o (Triangulabilidad de Superficies)}
\label{sec:app:rado}
%=======================================================================

El teorema de Rad\'o establece que toda superficie compacta admite una
triangulaci\'on, resultado fundamental tanto para la demostraci\'on del teorema de
Gauss--Bonnet (\Cref{thm:app:gauss-bonnet-general}) como para las aplicaciones del
aprendizaje profundo geom\'etrico sobre mallas trianguladas.

\subsection{Preliminares}
\label{subsec:app:rado-prelim}

Recordamos que una \emph{$2$-variedad topol\'ogica} (o \emph{superficie}) es un
espacio topol\'ogico $M$ que es Hausdorff, segundo contable, y tal que todo punto
$p \in M$ posee un entorno homeomorfo a un abierto de $\R^2$. El siguiente
resultado cl\'asico de topolog\'ia plana es una herramienta clave en la
demostraci\'on.

\begin{theorem}[Teorema de Jordan--Schoenflies]
\label{thm:app:jordan-schoenflies}
Todo homeomorfismo $h \colon S^1 \to \R^2$ cuya imagen es una curva de Jordan
$\mathcal{C}$ se extiende a un homeomorfismo
$H \colon \overline{D}^2 \to \overline{U}$, donde $D^2$ es el disco unitario
abierto y $U$ es la componente acotada de $\R^2 \setminus \mathcal{C}$.
\end{theorem}

Para una demostraci\'on, v\'ease \citet[Cap\'itulo~4]{moise1977geometric}.

\subsection{Enunciado formal}
\label{subsec:app:rado-statement}

\begin{theorem}[Teorema de Rad\'o, 1925]
\label{thm:app:rado}
Toda $2$-variedad topol\'ogica compacta admite una triangulaci\'on. Es decir, si
$M$ es una superficie compacta, existen un complejo simplicial $\mathcal{K}$ y un
homeomorfismo $h \colon \abs{\mathcal{K}} \to M$.
\end{theorem}

\begin{remark}[Triangulabilidad en dimensiones superiores]
\label{rem:app:triangulability-dimensions}
La triangulabilidad de variedades depende crucialmente de la dimensi\'on:
\begin{itemize}
    \item \textbf{Dimensi\'on 2}: Toda superficie compacta es triangulable
          (\Cref{thm:app:rado}).
    \item \textbf{Dimensi\'on 3}: \citet{moise1977geometric} demostr\'o en 1952
          que toda $3$-variedad compacta admite una triangulaci\'on, un resultado
          considerablemente m\'as dif\'icil.
    \item \textbf{Dimensi\'on $\geq 4$}: Existen variedades topol\'ogicas
          compactas no triangulables. El resultado de Casson (basado en el trabajo
          de Freedman sobre $4$-variedades) muestra que la variedad $E_8$ no
          admite triangulaci\'on alguna.
\end{itemize}
\end{remark}

\subsection{Esquema de demostraci\'on}
\label{subsec:app:rado-proof}

Presentamos las ideas principales de la prueba, siguiendo la exposici\'on de
\citet[Cap\'itulo~8]{moise1977geometric}. V\'ease tambi\'en
\citet[\S77]{munkres2000topology} para un tratamiento alternativo.

\begin{proof}[Esquema de demostraci\'on del \Cref{thm:app:rado}]
La demostraci\'on procede en cuatro pasos.

\textbf{Paso 1: Triangulabilidad local.}
Como $M$ es una $2$-variedad, cada punto $p \in M$ posee un entorno $U_p$
homeomorfo a un abierto de $\R^2$ mediante una carta $(U_p, \phi_p)$. En $\R^2$,
toda regi\'on acotada admite una triangulaci\'on (por ejemplo, mediante
subdivisi\'on baric\'entrica de una rejilla). Por tanto, cada $U_p$ puede
triangularse localmente transportando esta triangulaci\'on v\'ia $\phi_p^{-1}$.

\textbf{Paso 2: Cubrimiento finito.}
Por la compacidad de $M$, existe un subrecubrimiento finito $\{U_1, \ldots, U_N\}$
con triangulaciones locales $\mathcal{T}_1, \ldots, \mathcal{T}_N$.

\textbf{Paso 3: Compatibilidad mediante Jordan--Schoenflies.}
El paso esencial consiste en hacer compatibles las triangulaciones locales en las
regiones de solapamiento $U_i \cap U_j$. El \Cref{thm:app:jordan-schoenflies}
permite extender homeomorfismos entre curvas cerradas simples en $\R^2$ a
homeomorfismos de sus interiores, lo que garantiza que las triangulaciones locales
pueden ajustarse para ser compatibles en las intersecciones.

\textbf{Paso 4: Refinamiento y pegado.}
Mediante subdivisiones sucesivas (refinamientos) de las triangulaciones locales,
se obtienen triangulaciones suficientemente finas para ser compatibles en todas
las intersecciones. El pegado produce un complejo simplicial global $\mathcal{K}$
cuya realizaci\'on geom\'etrica $\abs{\mathcal{K}}$ es homeomorfa a $M$.
\end{proof}

%=======================================================================
\section{El Teorema Espectral para Matrices Sim\'etricas Reales}
\label{sec:app:spectral}
%=======================================================================

El teorema espectral es un resultado fundamental del \'algebra lineal que
garantiza la diagonalizaci\'on ortogonal de toda matriz sim\'etrica real. En el
contexto del aprendizaje profundo geom\'etrico, este teorema sustenta la
descomposici\'on espectral del laplaciano de grafo, la transformada de Fourier en
grafos y la definici\'on de matrices semidefinidas positivas. Para una exposici\'on
detallada, v\'ease \citet[Cap\'itulo~7]{Axler2015}.

\begin{theorem}[Teorema espectral para matrices sim\'etricas reales]
\label{thm:app:spectral}
Sea $A \in \R^{n \times n}$ una matriz sim\'etrica ($A = A^T$). Entonces:
\begin{enumerate}
    \item Todos los valores propios de $A$ son reales.
    \item Los vectores propios correspondientes a valores propios distintos son
          ortogonales.
    \item Existe una matriz ortogonal $U \in \Orth(n)$ tal que
          \[
          A = U \Lambda U^T , \quad
          \Lambda = \diag(\lambda_1, \ldots, \lambda_n) ,
          \]
          donde $\lambda_1, \ldots, \lambda_n$ son los valores propios de $A$.
\end{enumerate}
\end{theorem}

\begin{proof}
Procedemos por inducci\'on sobre $n$.

\textbf{Caso base $n = 1$.} Toda matriz $A = (a) \in \R^{1 \times 1}$ tiene un
\'unico valor propio $\lambda_1 = a \in \R$, y $U = (1)$ satisface
$A = U \Lambda U^T$ trivialmente.

\textbf{Paso inductivo.} Supongamos el resultado v\'alido para matrices
sim\'etricas de tama\~no $(n-1) \times (n-1)$ y sea $A \in \R^{n \times n}$
sim\'etrica.

\textbf{(1) Valores propios reales.} Sea $\lambda$ un valor propio de $A$
(posiblemente complejo) con vector propio
$\mathbf{v} \in \C^n \setminus \{\mathbf{0}\}$. Usando el producto herm\'itico
est\'andar $\inner{\mathbf{u}}{\mathbf{w}}_{\C} = \overline{\mathbf{u}}^T
\mathbf{w}$,
\[
\lambda \inner{\mathbf{v}}{\mathbf{v}}_{\C}
= \inner{\mathbf{v}}{A\mathbf{v}}_{\C}
= \inner{A^T \mathbf{v}}{\mathbf{v}}_{\C}
= \inner{A\mathbf{v}}{\mathbf{v}}_{\C}
= \overline{\inner{\mathbf{v}}{A\mathbf{v}}_{\C}}
= \bar{\lambda} \inner{\mathbf{v}}{\mathbf{v}}_{\C} ,
\]
donde usamos que $A = A^T$ implica $A = \bar{A}^T = A^*$ (pues $A$ es real). Como
$\inner{\mathbf{v}}{\mathbf{v}}_{\C} = \norm{\mathbf{v}}^2 > 0$, se obtiene
$\lambda = \bar{\lambda}$, es decir, $\lambda \in \R$.

\textbf{(2) Ortogonalidad de los vectores propios.} Sean
$A\mathbf{v} = \lambda \mathbf{v}$ y $A\mathbf{w} = \mu \mathbf{w}$ con
$\lambda \neq \mu$ (ambos reales por (1), y $\mathbf{v}, \mathbf{w} \in \R^n$).
Entonces
\[
\lambda \inner{\mathbf{v}}{\mathbf{w}} = \inner{A\mathbf{v}}{\mathbf{w}}
= \inner{\mathbf{v}}{A\mathbf{w}} = \mu \inner{\mathbf{v}}{\mathbf{w}} ,
\]
donde $\inner{\cdot}{\cdot}$ es el producto interior euclidiano en $\R^n$ y la
segunda igualdad usa la simetr\'ia de $A$. Luego
$(\lambda - \mu)\inner{\mathbf{v}}{\mathbf{w}} = 0$, y como $\lambda \neq \mu$,
concluimos $\inner{\mathbf{v}}{\mathbf{w}} = 0$.

\textbf{(3) Diagonalizaci\'on ortogonal.} Por (1), $A$ posee un valor propio real
$\lambda_1$ con vector propio unitario $\mathbf{v}_1 \in \R^n$
($\norm{\mathbf{v}_1} = 1$). Consideremos el complemento ortogonal
$W = \mathbf{v}_1^\perp = \{\mathbf{w} \in \R^n : \inner{\mathbf{w}}{\mathbf{v}_1}
= 0\}$, un subespacio de dimensi\'on $n - 1$.

\textit{Afirmaci\'on}: $W$ es $A$-invariante, es decir, $A\mathbf{w} \in W$ para
todo $\mathbf{w} \in W$. En efecto, si $\mathbf{w} \in W$,
\[
\inner{A\mathbf{w}}{\mathbf{v}_1} = \inner{\mathbf{w}}{A\mathbf{v}_1}
= \lambda_1 \inner{\mathbf{w}}{\mathbf{v}_1} = 0 ,
\]
por lo que $A\mathbf{w} \in \mathbf{v}_1^\perp = W$.

Sea $\{\mathbf{e}_2, \ldots, \mathbf{e}_n\}$ una base ortonormal de $W$ y
$P = [\mathbf{e}_2 \mid \cdots \mid \mathbf{e}_n] \in \R^{n \times (n-1)}$. La
restricci\'on $A|_W$ queda representada por la matriz sim\'etrica
$B = P^T A P \in \R^{(n-1) \times (n-1)}$. Por la hip\'otesis de inducci\'on,
existe una matriz ortogonal $\tilde{U} \in \Orth(n-1)$ con
$B = \tilde{U} \tilde{\Lambda} \tilde{U}^T$, donde
$\tilde{\Lambda} = \diag(\lambda_2, \ldots, \lambda_n)$. Definimos
\[
U = \begin{pmatrix} \mathbf{v}_1 & P\tilde{U} \end{pmatrix} \in \R^{n \times n} .
\]
Entonces $U$ es ortogonal (sus columnas forman una base ortonormal de $\R^n$) y
\[
U^T A U = \diag(\lambda_1, \lambda_2, \ldots, \lambda_n) ,
\]
lo que equivale a $A = U \Lambda U^T$, completando la inducci\'on.
\end{proof}

\begin{corollary}[Caracterizaci\'on de matrices semidefinidas positivas]
\label{cor:app:psd-characterization}
Sea $A \in \R^{n \times n}$ sim\'etrica con valores propios
$\lambda_1, \ldots, \lambda_n$. Entonces $A$ es semidefinida positiva
($A \succeq 0$) si y solo si $\lambda_i \geq 0$ para todo
$i \in \{1, \ldots, n\}$.
\end{corollary}

\begin{proof}
Por el \Cref{thm:app:spectral}, $A = U \Lambda U^T$ con $U$ ortogonal y
$\Lambda = \diag(\lambda_1, \ldots, \lambda_n)$. Para cualquier
$\mathbf{x} \in \R^n$, sea $\mathbf{y} = U^T \mathbf{x}$. Entonces
\[
\mathbf{x}^T A \mathbf{x} = \mathbf{x}^T U \Lambda U^T \mathbf{x}
= \mathbf{y}^T \Lambda \mathbf{y} = \sum_{i=1}^n \lambda_i y_i^2 .
\]
Como $U$ es ortogonal, la aplicaci\'on $\mathbf{x} \mapsto \mathbf{y} = U^T
\mathbf{x}$ es una biyecci\'on de $\R^n$. Si todos los $\lambda_i \geq 0$, la suma
es claramente $\geq 0$ para todo $\mathbf{x}$. Rec\'iprocamente, si alg\'un
$\lambda_j < 0$, tomando $\mathbf{x} = U\mathbf{e}_j$ (donde $\mathbf{e}_j$ es el
$j$-\'esimo vector can\'onico) se obtiene
$\mathbf{x}^T A \mathbf{x} = \lambda_j < 0$.
\end{proof}

%=======================================================================
\section{El Teorema de la Funci\'on Inversa}
\label{sec:app:inverse-function}
%=======================================================================

El teorema de la funci\'on inversa es una herramienta fundamental del an\'alisis
que garantiza la existencia local de inversas diferenciables. En el contexto de
las variedades diferenciables se usa repetidamente para construir difeomorfismos
locales, en particular para demostrar que el mapa exponencial riemanniano es un
difeomorfismo local.

\subsection{Preliminar: el teorema del punto fijo de Banach}
\label{subsec:app:banach-fixed-point}

\begin{theorem}[Teorema del punto fijo de Banach]
\label{thm:app:banach-fixed-point}
Sea $(X, d)$ un espacio m\'etrico completo y $T \colon X \to X$ una
\emph{contracci\'on}, es decir, existe $\lambda \in [0, 1)$ tal que
\[
d(T(x), T(y)) \leq \lambda \, d(x, y) \quad \text{para todo } x, y \in X .
\]
Entonces $T$ posee un \'unico punto fijo $x^* \in X$ (es decir, $T(x^*) = x^*$).
Adem\'as, para cualquier $x_0 \in X$, la sucesi\'on $x_{n+1} = T(x_n)$ converge a
$x^*$, con la estimaci\'on
\[
d(x_n, x^*) \leq \frac{\lambda^n}{1 - \lambda} \, d(x_0, x_1) .
\]
\end{theorem}

\begin{proof}
\textbf{Existencia.} Sea $x_0 \in X$ arbitrario y definamos $x_{n+1} = T(x_n)$.
Para $n \geq 1$,
\[
d(x_{n+1}, x_n) = d(T(x_n), T(x_{n-1})) \leq \lambda \, d(x_n, x_{n-1})
\leq \cdots \leq \lambda^n \, d(x_1, x_0) .
\]
Para $m > n$, la desigualdad triangular da
\[
d(x_m, x_n) \leq \sum_{k=n}^{m-1} d(x_{k+1}, x_k)
\leq \sum_{k=n}^{m-1} \lambda^k \, d(x_1, x_0)
\leq \frac{\lambda^n}{1 - \lambda} \, d(x_1, x_0) .
\]
Como $\lambda < 1$, el lado derecho tiende a $0$ cuando $n \to \infty$, luego
$(x_n)$ es una sucesi\'on de Cauchy. Por la completitud de $(X, d)$, existe el
l\'imite $x^* = \lim_{n \to \infty} x_n$.

Por la continuidad de $T$ (toda contracci\'on es lipschitziana, luego continua),
\[
T(x^*) = T\!\left(\lim_{n \to \infty} x_n\right)
= \lim_{n \to \infty} T(x_n) = \lim_{n \to \infty} x_{n+1} = x^* .
\]

\textbf{Unicidad.} Si $T(y^*) = y^*$ para otro punto $y^* \in X$, entonces
\[
d(x^*, y^*) = d(T(x^*), T(y^*)) \leq \lambda \, d(x^*, y^*) .
\]
Como $\lambda < 1$, esto fuerza $d(x^*, y^*) = 0$, es decir, $x^* = y^*$.

\textbf{Estimaci\'on.} Tomando $m \to \infty$ en la desigualdad
$d(x_m, x_n) \leq \frac{\lambda^n}{1-\lambda} d(x_1, x_0)$ se obtiene el resultado.
\end{proof}

\subsection{Enunciado y demostraci\'on del teorema de la funci\'on inversa}
\label{subsec:app:inverse-function-proof}

\begin{theorem}[Teorema de la funci\'on inversa]
\label{thm:app:inverse-function}
Sea $U \subseteq \R^n$ un abierto, $F \colon U \to \R^n$ una aplicaci\'on de clase
$C^k$ ($k \geq 1$), y $p \in U$ un punto tal que el diferencial
$DF_p \colon \R^n \to \R^n$ es invertible. Entonces existen entornos abiertos $V$
de $p$ y $W$ de $F(p)$ tales que $F|_V \colon V \to W$ es un difeomorfismo de
clase $C^k$.
\end{theorem}

\begin{proof}
La demostraci\'on sigue la estrategia cl\'asica de
\citet[Teorema~9.24]{Rudin1976}, que presentamos en cuatro pasos.

\textbf{Paso 1: Reducci\'on al caso $DF_p = I$.}
Consideramos $G = (DF_p)^{-1} \circ F$. Entonces $G$ es de clase $C^k$,
$G(p) = (DF_p)^{-1}(F(p))$ y $DG_p = (DF_p)^{-1} \circ DF_p = I$. Si demostramos el
teorema para $G$, la composici\'on con la aplicaci\'on lineal invertible $DF_p$
(un difeomorfismo global) proporciona el resultado para $F$. As\'i, podemos
suponer sin p\'erdida de generalidad que $DF_p = I$ y, tras trasladar, que $p = 0$
y $F(0) = 0$.

\textbf{Paso 2: Sobreyectividad local v\'ia el punto fijo de Banach.}
Definimos $\varphi(x) = x - F(x)$, de modo que $D\varphi_0 = I - DF_0 = 0$. Por la
continuidad de $D\varphi$, existe $r > 0$ con
$\norm{D\varphi_x} \leq \tfrac{1}{2}$ para todo $x \in \overline{B}(0, r)$. Por el
teorema del valor medio, para $x_1, x_2 \in \overline{B}(0, r)$,
\[
\norm{\varphi(x_1) - \varphi(x_2)} \leq \tfrac{1}{2} \norm{x_1 - x_2} .
\]
Para cada $y \in B(0, r/2)$, definimos
$T_y \colon \overline{B}(0, r) \to \R^n$ por
$T_y(x) = y + \varphi(x) = y + x - F(x)$. Entonces
\[
\norm{T_y(x)} \leq \norm{y} + \norm{\varphi(x) - \varphi(0)}
\leq \tfrac{r}{2} + \tfrac{1}{2}\norm{x} \leq r ,
\]
luego $T_y$ aplica $\overline{B}(0, r)$ en s\'i misma. Adem\'as, $T_y$ es una
contracci\'on con constante $\tfrac{1}{2}$, pues
$T_y(x_1) - T_y(x_2) = \varphi(x_1) - \varphi(x_2)$. Como $\overline{B}(0, r)$ es
completo, el \Cref{thm:app:banach-fixed-point} garantiza un \'unico
$x \in \overline{B}(0, r)$ con $T_y(x) = x$, equivalentemente $F(x) = y$. Esto
muestra que $F$ es sobreyectiva sobre $W = B(0, r/2)$.

\textbf{Paso 3: Inyectividad.}
Si $F(x_1) = F(x_2)$ con $x_1, x_2 \in \overline{B}(0, r)$, entonces
\[
\norm{x_1 - x_2} = \norm{\varphi(x_1) - \varphi(x_2) + F(x_1) - F(x_2)}
= \norm{\varphi(x_1) - \varphi(x_2)} \leq \tfrac{1}{2}\norm{x_1 - x_2} ,
\]
de donde $x_1 = x_2$. Sea $V = F^{-1}(W) \cap B(0, r)$; entonces
$F|_V \colon V \to W$ es biyectiva.

\textbf{Paso 4: Diferenciabilidad de $F^{-1}$.}
Sea $g = (F|_V)^{-1} \colon W \to V$. Debemos mostrar que $g$ es de clase $C^k$.
Para $y, y + k \in W$ con $x = g(y)$ y $x + h = g(y + k)$,
\[
k = F(x + h) - F(x) = DF_x(h) + \norm{h} \cdot \varepsilon(h) ,
\]
donde $\varepsilon(h) \to 0$ cuando $h \to 0$. De la estimaci\'on de
contracci\'on, $\norm{h} \leq 2\norm{k}$, por lo que $h \to 0$ cuando $k \to 0$.
Como $DF_x$ es invertible (su determinante es una funci\'on continua no nula),
\[
g(y + k) - g(y) = h = (DF_x)^{-1}(k) + \norm{k} \cdot \tilde{\varepsilon}(k) ,
\]
con $\tilde{\varepsilon}(k) \to 0$, luego $Dg_y = (DF_{g(y)})^{-1}$. Como la
inversi\'on de matrices y la composici\'on son operaciones $C^\infty$, un
argumento inductivo muestra que $g$ es de clase $C^k$.
\end{proof}

%=======================================================================
\section{Preliminares para el Teorema de Aproximaci\'on Universal}
\label{sec:app:uat-preliminaries}
%=======================================================================

El teorema de aproximaci\'on universal del \Cref{ch:foundations} garantiza que
cualquier funci\'on continua puede aproximarse arbitrariamente bien por una red
neuronal artificial con una \'unica capa oculta. Su demostraci\'on requiere
herramientas de teor\'ia de la medida y an\'alisis funcional: las nociones de
medida y $\sigma$-\'algebra fundamentan los teoremas de representaci\'on (Riesz) y
separaci\'on (Hahn--Banach) que constituyen el n\'ucleo de la prueba. Reunimos
aqu\'i los preliminares necesarios, dejando el enunciado y la demostraci\'on del
teorema en s\'i para el \Cref{ch:foundations}. Para un tratamiento en profundidad
de estos fundamentos, v\'ease \citet{Folland1999} o \citet{Rudin1976}.

\subsection{Preliminares de teor\'ia de la medida}
\label{subsec:app:measure-theory}

\begin{definition}[$\sigma$-\'algebra]
\label{def:app:sigma-algebra}
Sea $X$ un conjunto no vac\'io. Una \emph{$\sigma$-\'algebra} (o $\sigma$-campo)
sobre $X$ es una colecci\'on $\mathcal{F}$ de subconjuntos de $X$ que satisface:
\begin{enumerate}
    \item $X \in \mathcal{F}$;
    \item si $A \in \mathcal{F}$, entonces $A^c = X \setminus A \in \mathcal{F}$
          (cerrada bajo complementos);
    \item si $\{A_n\}_{n=1}^\infty \subseteq \mathcal{F}$, entonces
          $\bigcup_{n=1}^\infty A_n \in \mathcal{F}$ (cerrada bajo uniones
          numerables).
\end{enumerate}
Los elementos de $\mathcal{F}$ se denominan \emph{conjuntos medibles}.
\end{definition}

\begin{definition}[Espacio medible]
\label{def:app:measurable-space}
Un \emph{espacio medible} es un par $(X, \mathcal{F})$ donde $X$ es un conjunto no
vac\'io y $\mathcal{F}$ es una $\sigma$-\'algebra sobre $X$.
\end{definition}

\begin{definition}[$\sigma$-\'algebra de Borel]
\label{def:app:borel-sigma-algebra}
Sea $(X, \tau)$ un espacio topol\'ogico. La \emph{$\sigma$-\'algebra de Borel}
$\mathcal{B}(X)$ es la $\sigma$-\'algebra m\'as peque\~na que contiene todos los
conjuntos abiertos de $X$. Sus elementos se denominan \emph{conjuntos de Borel} o
\emph{borelianos}.
\end{definition}

En particular, para $\R^n$ la $\sigma$-\'algebra de Borel $\mathcal{B}(\R^n)$
est\'a generada por los rect\'angulos abiertos, o equivalentemente por los
semiespacios de la forma $\{x \in \R^n : w^T x + b > 0\}$.

\begin{definition}[Medida]
\label{def:app:measure}
Sea $(X, \mathcal{F})$ un espacio medible. Una \emph{medida} es una funci\'on
$\mu \colon \mathcal{F} \to [0, +\infty]$ que satisface:
\begin{enumerate}
    \item $\mu(\emptyset) = 0$;
    \item \textbf{$\sigma$-aditividad}: para cualquier sucesi\'on
          $\{A_n\}_{n=1}^\infty$ de conjuntos disjuntos dos a dos en
          $\mathcal{F}$,
          \[
          \mu\!\left(\bigcup_{n=1}^\infty A_n\right) = \sum_{n=1}^\infty \mu(A_n) .
          \]
\end{enumerate}
\end{definition}

\begin{definition}[Espacio de medida]
\label{def:app:measure-space}
Un \emph{espacio de medida} es una terna $(X, \mathcal{F}, \mu)$ donde $X$ es un
conjunto no vac\'io, $\mathcal{F}$ es una $\sigma$-\'algebra sobre $X$ y
$\mu \colon \mathcal{F} \to [0, +\infty]$ es una medida. Si $\mu(X) < \infty$,
decimos que $\mu$ es una \emph{medida finita}. Si $\mu(X) = 1$, decimos que $\mu$
es una \emph{medida de probabilidad}.
\end{definition}

\begin{definition}[Medida finita con signo]
\label{def:app:signed-measure}
Sea $(X, \mathcal{F})$ un espacio medible. Una \emph{medida finita con signo} es
una funci\'on $\mu \colon \mathcal{F} \to \R$ con $\mu(\emptyset) = 0$ que es
$\sigma$-aditiva: para cualquier sucesi\'on $\{A_n\}_{n=1}^\infty$ de conjuntos
disjuntos dos a dos en $\mathcal{F}$,
\[
\mu\!\left(\bigcup_{n=1}^\infty A_n\right) = \sum_{n=1}^\infty \mu(A_n) ,
\]
donde la serie converge absolutamente. La \emph{variaci\'on total} de $\mu$ es la
medida $\abs{\mu}$ definida por
\[
\abs{\mu}(A) = \sup\!\left\{ \sum_{i=1}^n \abs{\mu(A_i)} :
\{A_i\}_{i=1}^n \text{ partici\'on finita de } A \right\} .
\]
\end{definition}

\begin{definition}[Medida de Lebesgue]
\label{def:app:lebesgue-measure}
La \emph{medida de Lebesgue} $\lambda$ en $\R^n$ es la \'unica medida en
$(\R^n, \mathcal{B}(\R^n))$ que extiende la noci\'on de volumen $n$-dimensional y
satisface:
\begin{itemize}
    \item para un rect\'angulo $R = [a_1, b_1] \times \cdots \times [a_n, b_n]$,
          $\lambda(R) = \prod_{i=1}^n (b_i - a_i)$;
    \item es invariante bajo traslaciones: $\lambda(A + x) = \lambda(A)$ para todo
          $x \in \R^n$.
\end{itemize}
\end{definition}

Un conjunto $A$ tiene \emph{medida de Lebesgue nula} (o medida cero) si
$\lambda(A) = 0$. Ejemplos incluyen puntos aislados, conjuntos finitos, conjuntos
numerables como $\Q$ y el conjunto de Cantor.

\begin{definition}[Funci\'on medible]
\label{def:app:measurable-function}
Sean $(X, \mathcal{F})$ y $(Y, \mathcal{G})$ espacios medibles. Una funci\'on
$f \colon X \to Y$ es \emph{medible} (o $\mathcal{F}$-$\mathcal{G}$-medible) si la
preimagen de todo conjunto medible es medible:
\[
f^{-1}(B) = \{x \in X : f(x) \in B\} \in \mathcal{F}
\quad \forall B \in \mathcal{G} .
\]
\end{definition}

Cuando $Y = \R$ con la $\sigma$-\'algebra de Borel, basta verificar que
$\{x : f(x) > \alpha\} \in \mathcal{F}$ para todo $\alpha \in \R$. En particular,
toda funci\'on continua es medible respecto a las $\sigma$-\'algebras de Borel.

\begin{definition}[Integral de Lebesgue]
\label{def:app:lebesgue-integral}
Sea $(X, \mathcal{F}, \mu)$ un espacio de medida. La \emph{integral de Lebesgue}
se construye en tres pasos:
\begin{enumerate}
    \item \textbf{Funciones simples}: para $s = \sum_{i=1}^n a_i \chi_{A_i}$ con
          $a_i \geq 0$ y $A_i \in \mathcal{F}$ disjuntos, donde $\chi_A$ es la
          funci\'on caracter\'istica de $A$, se define
          \[
          \int_X s \, d\mu = \sum_{i=1}^n a_i \mu(A_i) .
          \]
    \item \textbf{Funciones no negativas}: para $f \colon X \to [0, +\infty]$
          medible,
          \[
          \int_X f \, d\mu = \sup\!\left\{\int_X s \, d\mu :
          s \text{ simple, } 0 \leq s \leq f\right\} .
          \]
    \item \textbf{Funciones generales}: para $f \colon X \to \R$ medible,
          escribimos $f = f^+ - f^-$ con $f^+ = \max(f, 0)$ y $f^- = \max(-f, 0)$.
          Decimos que $f$ es \emph{integrable} (o $f \in L^1(\mu)$) si
          $\int_X f^+ \, d\mu < \infty$ y $\int_X f^- \, d\mu < \infty$, y se
          define
          \[
          \int_X f \, d\mu = \int_X f^+ \, d\mu - \int_X f^- \, d\mu .
          \]
\end{enumerate}
\end{definition}

\begin{definition}[Casi por doquier]
\label{def:app:almost-everywhere}
Sea $(X, \mathcal{F}, \mu)$ un espacio de medida. Una propiedad $P(x)$ se cumple
\emph{casi por doquier} (c.p.d.) si el conjunto de puntos donde no se cumple tiene
medida cero:
\[
\mu(\{x \in X : \neg P(x)\}) = 0 .
\]
\end{definition}

En el contexto de la medida de Lebesgue en $\R^n$, ``casi por doquier'' significa
``excepto en un conjunto de medida de Lebesgue nula''.

\subsection{Preliminares de an\'alisis funcional}
\label{subsec:app:functional-analysis}

\begin{definition}[Espacio vectorial]
\label{def:app:vector-space}
Un \emph{espacio vectorial} sobre $\R$ es un conjunto $V$ equipado con dos
operaciones, suma $+ \colon V \times V \to V$ y producto por escalares
$\cdot \colon \R \times V \to V$, que satisfacen los siguientes axiomas para todo
$u, v, w \in V$ y $\alpha, \beta \in \R$:
\begin{enumerate}
    \item \textbf{Asociatividad de la suma:} $(u + v) + w = u + (v + w)$;
    \item \textbf{Conmutatividad de la suma:} $u + v = v + u$;
    \item \textbf{Elemento neutro:} existe $0 \in V$ con $v + 0 = v$;
    \item \textbf{Elemento opuesto:} para cada $v$ existe $-v \in V$ con
          $v + (-v) = 0$;
    \item \textbf{Compatibilidad del producto:}
          $\alpha(\beta v) = (\alpha\beta)v$;
    \item \textbf{Elemento unidad:} $1 \cdot v = v$;
    \item \textbf{Distributividad respecto a la suma vectorial:}
          $\alpha(u + v) = \alpha u + \alpha v$;
    \item \textbf{Distributividad respecto a la suma escalar:}
          $(\alpha + \beta)v = \alpha v + \beta v$.
\end{enumerate}
\end{definition}

\begin{definition}[Espacio vectorial normado]
\label{def:app:normed-space}
Un \emph{espacio vectorial normado} es un par $(X, \norm{\cdot})$ donde $X$ es un
espacio vectorial real y $\norm{\cdot} \colon X \to [0, +\infty)$ es una
\emph{norma}, es decir, una funci\'on que satisface:
\begin{enumerate}
    \item $\norm{x} = 0$ si y solo si $x = 0$;
    \item $\norm{\alpha x} = \abs{\alpha} \norm{x}$ para todo $\alpha \in \R$;
    \item $\norm{x + y} \leq \norm{x} + \norm{y}$ (desigualdad triangular).
\end{enumerate}
\end{definition}

Ejemplos fundamentales incluyen $\R^n$ con la norma euclidiana y los espacios de
funciones $L^p(\Omega)$.

\begin{definition}[Recubrimiento]
\label{def:app:cover}
Sea $X$ un espacio topol\'ogico y $K \subseteq X$. Un \emph{recubrimiento} de $K$
es una familia de conjuntos $\{U_\alpha\}_{\alpha \in A}$ con
$K \subseteq \bigcup_{\alpha \in A} U_\alpha$. Si cada $U_\alpha$ es abierto, es un
\emph{recubrimiento abierto}. Un \emph{subrecubrimiento} es una subfamilia que
sigue siendo un recubrimiento de $K$.
\end{definition}

\begin{definition}[Espacio compacto]
\label{def:app:compact-space}
Un subconjunto $K$ de un espacio topol\'ogico es \emph{compacto} si todo
recubrimiento abierto de $K$ admite un subrecubrimiento finito.
\end{definition}

Intuitivamente, la compacidad generaliza las propiedades de los conjuntos cerrados
y acotados en $\R^n$ a espacios topol\'ogicos generales.

\begin{theorem}[Heine--Borel]
\label{thm:app:heine-borel}
Un subconjunto $K \subset \R^n$ es compacto si y solo si es cerrado y acotado
\citep[Teorema~2.41]{Rudin1976}.
\end{theorem}

En particular, el hipercubo unitario $I_m = [0,1]^m$ es compacto, lo cual es
esencial para el teorema de aproximaci\'on universal.

\begin{definition}[Funcional lineal]
\label{def:app:linear-functional}
Sea $X$ un espacio vectorial real. Un \emph{funcional lineal} es una aplicaci\'on
$f \colon X \to \R$ que satisface
\[
f(\alpha x + \beta y) = \alpha f(x) + \beta f(y)
\quad \forall x, y \in X, \, \forall \alpha, \beta \in \R .
\]
Es decir, un funcional lineal es un elemento del espacio dual algebraico $X^*$.
\end{definition}

\begin{definition}[Funcional lineal continuo]
\label{def:app:continuous-functional}
Sea $(X, \norm{\cdot})$ un espacio vectorial normado. Un funcional lineal
$f \colon X \to \R$ es \emph{continuo} (o \emph{acotado}) si existe una constante
$M > 0$ con
\[
\abs{f(x)} \leq M\norm{x} \quad \forall x \in X .
\]
\end{definition}

El conjunto de todos los funcionales lineales continuos en $X$ forma el
\emph{espacio dual} $X'$, un espacio de Banach con la norma del operador
$\norm{f} = \sup_{\norm{x} \leq 1} \abs{f(x)}$.

\begin{definition}[Espacio $C(K)$]
\label{def:app:space-CK}
Sea $K$ un espacio compacto. El espacio $C(K)$ es el conjunto de todas las
funciones continuas $f \colon K \to \R$ equipado con la \emph{norma supremo} (o
norma uniforme)
\[
\norm{f}_\infty = \sup_{x \in K} \abs{f(x)} .
\]
\end{definition}

Con esta norma, $C(K)$ es un espacio de Banach (un espacio vectorial normado
completo). En el contexto del teorema de aproximaci\'on universal, trabajamos con
$C(I_m)$, el espacio de funciones continuas sobre el hipercubo unitario.

\begin{theorem}[Convergencia dominada de Lebesgue]
\label{thm:app:dominated-convergence}
Sea $(X, \mathcal{F}, \mu)$ un espacio de medida y $\{f_n\}_{n=1}^\infty$ una
sucesi\'on de funciones medibles $f_n \colon X \to \R$ tal que:
\begin{enumerate}
    \item $f_n(x) \to f(x)$ para casi todo $x \in X$;
    \item existe una funci\'on integrable $g \colon X \to \R$ con
          $\abs{f_n(x)} \leq g(x)$ para casi todo $x \in X$ y todo
          $n \in \mathbb{N}$.
\end{enumerate}
Entonces $f$ es integrable y
\[
\lim_{n \to \infty} \int_X f_n \, d\mu = \int_X f \, d\mu
= \int_X \lim_{n \to \infty} f_n \, d\mu
\]
\citep[Teorema~2.24]{Folland1999}.
\end{theorem}

Este teorema es fundamental para intercambiar l\'imites e integrales, lo cual es
necesario para demostrar que las funciones sigmoidales son discriminatorias.

\begin{definition}[Conjunto convexo]
\label{def:app:convex-set}
Un subconjunto $C$ de un espacio vectorial $V$ es \emph{convexo} si para todo
$x, y \in C$ y $\lambda \in [0,1]$ se tiene $\lambda x + (1-\lambda)y \in C$.
Equivalentemente, $C$ contiene el segmento que une cualquier par de sus puntos.
\end{definition}

\begin{theorem}[Hahn--Banach, forma geom\'etrica]
\label{thm:app:hahn-banach}
Sea $X$ un espacio vectorial topol\'ogico real, y $A$, $B$ subconjuntos convexos
disjuntos de $X$ con $A$ abierto. Entonces existen un funcional lineal continuo
$f \colon X \to \R$ y una constante $\alpha \in \R$ tales que
\[
f(a) < \alpha \leq f(b) \quad \forall a \in A, \, b \in B
\]
\citep[Teorema~5.5]{Folland1999}.
\end{theorem}

\begin{corollary}[Caracterizaci\'on de densidad]
\label{cor:app:density-characterization}
Sea $X$ un espacio vectorial normado y $Y \subset X$ un subespacio. Entonces $Y$
es denso en $X$ si y solo si el \'unico funcional lineal continuo que se anula en
$Y$ es el funcional cero.
\end{corollary}

\begin{proof}
\textbf{($\Rightarrow$)} Supongamos que $Y$ es denso en $X$. Sea
$f \colon X \to \R$ un funcional lineal continuo con $f(y) = 0$ para todo
$y \in Y$. Para cualquier $x \in X$ existe una sucesi\'on $\{y_n\} \subset Y$ con
$y_n \to x$. Por la continuidad de $f$,
\[
f(x) = f\!\left(\lim_{n \to \infty} y_n\right)
= \lim_{n \to \infty} f(y_n) = 0 ,
\]
por lo que $f \equiv 0$.

\textbf{($\Leftarrow$)} Supongamos que $Y$ no es denso en $X$. Entonces
$\overline{Y} \neq X$, y existe $x_0 \in X \setminus \overline{Y}$. Como
$\overline{Y}$ es un subespacio cerrado, tomamos una bola abierta
$B_r(0) \subset X$ y definimos $A = \overline{Y} + B_r(0)$, que es abierto y
convexo con $A \cap \{x_0\} = \emptyset$ para $r$ suficientemente peque\~no (ya que
$x_0 \notin \overline{Y}$). Por el \Cref{thm:app:hahn-banach}, existen un funcional
lineal $f$ y $\alpha \in \R$ con
\[
f(a) < \alpha \leq f(x_0) \quad \forall a \in A .
\]
Como $\overline{Y}$ es un subespacio, $0 \in \overline{Y} \subset A$, luego
$f(0) = 0 < \alpha$. Adem\'as, para cualquier $y \in Y$ y $t \in \R$, $ty \in
\overline{Y} \subset A$, de modo que $tf(y) = f(ty) < \alpha$ para todo $t$. Esto
solo es posible si $f(y) = 0$. Sin embargo, $f(x_0) \geq \alpha > 0$, por lo que
$f \not\equiv 0$.
\end{proof}

\begin{theorem}[Representaci\'on de Riesz para $C(K)$]
\label{thm:app:riesz}
Sea $K$ un espacio compacto y $C(K)$ el espacio de funciones continuas reales en
$K$ con la norma supremo. Para todo funcional lineal continuo
$L \colon C(K) \to \R$ existe una \'unica medida de Borel regular finita con signo
$\mu$ en $K$ tal que
\[
L(f) = \int_K f \, d\mu \quad \forall f \in C(K) .
\]
Adem\'as, $\norm{L} = \abs{\mu}(K)$, donde $\abs{\mu}$ es la variaci\'on total de
$\mu$ \citep[Teorema~7.17]{Folland1999}.
\end{theorem}

\begin{theorem}[Stone--Weierstrass]
\label{thm:app:stone-weierstrass}
Sea $K$ un espacio compacto y $\mathcal{A} \subset C(K)$ una sub\'algebra que
separa puntos de $K$ y no se anula en ning\'un punto. Entonces $\mathcal{A}$ es
densa en $C(K)$ con respecto a la norma supremo
\citep[Teorema~7.32]{Rudin1976}.
\end{theorem}

En particular, si $\mathcal{A}$ contiene las funciones constantes y separa puntos,
entonces es densa en $C(K)$.

\begin{remark}
Existen dos enfoques principales para demostrar el teorema de aproximaci\'on
universal. El enfoque de \citet{cybenko1989} usa el corolario de Hahn--Banach
junto con el teorema de representaci\'on de Riesz, demostrando que las funciones
sigmoidales son ``discriminatorias'' (se anulan contra una medida solo si la
medida es cero). El enfoque de \citet{hornik1989} usa el teorema de
Stone--Weierstrass, demostrando que las redes neuronales forman una sub\'algebra
que separa puntos. Ambos se apoyan en los preliminares reunidos arriba, y el
enunciado y la demostraci\'on completos se desarrollan en el \Cref{ch:foundations}.
\end{remark}
